\documentclass[a4paper,12pt]{article}

% Пакеты для русского языка и кодировки
\usepackage[utf8]{inputenc}
\usepackage[T2A]{fontenc}
\usepackage[russian]{babel}

% Пакеты для математики
\usepackage{amsmath}
\usepackage{amssymb}
\usepackage{amsthm}

% Пакеты для оформления
\usepackage{geometry}
\geometry{left=2.5cm, right=2.5cm, top=2.5cm, bottom=2.5cm}
\usepackage{hyperref}
\usepackage{enumitem}
\usepackage{booktabs}
\usepackage{array}

% Настройка заголовков
\title{2.3.5 Материалы для изучения по научной специальности \\
Математическое и программное обеспечение вычислительных систем, комплексов и компьютерных сетей \\
Гусев Виталий Евгеньевич 2.3.5 2026}
\date{}

\begin{document}

\maketitle

\tableofcontents
\newpage

\setcounter{section}{2} % чтобы раздел имел номер 2
\section{Раздел 2. Дискретная математика и математическая логика.}

\subsection{Основы комбинаторного анализа. Метод производящих функций, метод включений и исключений. Примеры применения}

\noindent\textit{Основной источник: Липский В. Комбинаторика для программистов. М.: Мир, 1988. Гл. 1, §1.1--1.13, с. 9--76.}

\subsubsection{Основные понятия комбинаторики: размещения, перестановки, сочетания}

\noindent\textit{Основные источники: Липский В. Комбинаторика для программистов. М.: Мир, 1988. Гл. 1, §1.2--1.3, с. 14--24; §1.5--1.6, с. 34--45.}

\paragraph{Постановка комбинаторных задач о размещениях.}
Классическая задача комбинаторики состоит в подсчёте числа способов размещения объектов по «ящикам» при заданных ограничениях. Формально даны два конечных множества \(X\) и \(Y\), причём
\[
|X|=n,\qquad |Y|=m.
\]
Элементы \(X\) интерпретируются как объекты, элементы \(Y\) --- как ящики, а произвольная функция \(f:X\to Y\) задаёт размещение: объект \(x\in X\) попадает в ящик \(f(x)\in Y\). Без потери общности можно считать
\[
X=\{1,\dots,n\},\qquad Y=\{1,\dots,m\},
\]
отождествляя функцию \(f\) с последовательностью
\[
\langle f(1),\dots,f(n)\rangle.
\]
В зависимости от ограничений на \(f\) получаются различные комбинаторные объекты: произвольные функции, инъективные функции (размещения без повторений), перестановки, упорядоченные размещения по ящикам, сочетания и сочетания с повторениями.
\textit{(Источник: Липский, гл. 1, §1.2, с. 14--15.)}

\paragraph{Функции и размещения без ограничений.}
Если на размещение не накладывается никаких ограничений, каждый объект может быть помещён в любой из \(m\) ящиков независимо от остальных. Справедлива следующая теорема.

\textbf{Теорема 1.1 (о числе функций).} Если \(|X|=n\), \(|Y|=m\), то число всех функций \(f:X\to Y\) равно
\[
m^n.
\]

\emph{Доказательство.} Считая \(X=\{1,\dots,n\}\), сводим задачу к подсчёту числа всех последовательностей \(\langle y_1,\dots,y_n\rangle\) с членами из \(m\)-элементного множества \(Y\). Каждый член \(y_i\) можно выбрать \(m\) способами независимо от остальных. Поэтому число всех последовательностей равно
\[
\underbrace{m\cdot m\cdot\ldots\cdot m}_{n\ \text{раз}}=m^n.
\]
Теорема доказана.
\textit{(Источник: Липский, гл. 1, §1.2, с. 15.)}

\paragraph{Инъективные размещения и перестановки.}
Если каждый ящик содержит не более одного объекта, то различным объектам соответствуют различные ящики, то есть функция \(f:X\to Y\) должна быть взаимно однозначной (инъективной). Обозначим через \([m]_n\) число всех инъективных функций из \(n\)-элементного множества в \(m\)-элементное множество. Полагаем \([m]_0=1\).

\textbf{Теорема 1.2 (о числе инъективных функций).} Если \(|X|=n\), \(|Y|=m\), то число всех взаимно однозначных функций \(f:X\to Y\) равно
\[
[m]_n=m(m-1)\cdots(m-n+1).
\]
При \(n>m\) это число равно нулю.

\emph{Доказательство.} Будем строить инъективную последовательность \(\langle y_1,\dots,y_n\rangle\) с членами из \(Y\). Элемент \(y_1\) можно выбрать \(m\) способами. После выбора \(y_1,\dots,y_{i-1}\) элемент \(y_i\) можно выбрать любым из
\[
m-(i-1)=m-i+1
\]
элементов множества
\[
Y\setminus\{y_1,\dots,y_{i-1}\}.
\]
Перемножая возможности для всех \(i=1,\dots,n\), получаем
\[
m(m-1)\cdots(m-n+1)=[m]_n.
\]
Если \(n>m\), то на некотором шаге свободных ящиков не останется, поэтому инъективных функций нет.
\textit{(Источник: Липский, гл. 1, §1.2, с. 15.)}

В частности, если \(m=n\), то каждая инъективная функция \(f:X\to Y\) является взаимно однозначным отображением \(X\) на \(Y\). Такое отображение называется \emph{перестановкой} множества \(X\). Из теоремы 1.2 получаем:
\[
[n]_n=n(n-1)\cdots1=n!.
\]
Таким образом, число перестановок \(n\)-элементного множества равно \(n!\).
\textit{(Источник: Липский, гл. 1, §1.2, с. 16.)}

\textbf{Пример.} Число инъективных последовательностей длины \(3\) с элементами из множества \(\{1,2,3,4\}\) равно
\[
[4]_3=4\cdot3\cdot2=24.
\]
Это все упорядоченные тройки различных чисел из \(\{1,2,3,4\}\).
\textit{(Источник: Липский, гл. 1, §1.2, с. 16.)}

\paragraph{Упорядоченные размещения по ящикам.}
Рассмотрим теперь ситуацию, когда \(n\) объектов размещаются по \(m\) ящикам так, что каждый ящик содержит не множество, а \emph{последовательность} помещённых в него объектов. Два размещения считаются равными, если в каждом ящике содержится одна и та же последовательность объектов. Такие размещения называются \emph{упорядоченными размещениями} \(n\) объектов по \(m\) ящикам. Число таких размещений обозначается \([m]^n\); полагаем \([m]^0=1\).

\textbf{Теорема 1.4 (о числе упорядоченных размещений).} Число упорядоченных размещений \(n\) объектов по \(m\) ящикам равно
\[
[m]^n=m(m+1)\cdots(m+n-1).
\]

\emph{Доказательство.} Будем добавлять объекты по одному. Первый объект можно разместить \(m\) способами. Второй объект можно поместить либо в один из \(m-1\) пустых ящиков, либо в ящик, содержащий первый объект, --- перед ним или после него; всего \(m+1\) способ. В общем случае предположим, что уже размещено \(i-1\) объектов, причём в \(k\)-м ящике находятся \(r_k\) объектов. Тогда \(i\)-й объект можно добавить в \(k\)-й ящик \(r_k+1\) способами. Суммарное число возможностей равно
\[
(r_1+1)+\dots+(r_m+1)=(r_1+\dots+r_m)+m=(i-1)+m=m+i-1.
\]
Перемножая эти числа для \(i=1,\dots,n\), получаем
\[
m(m+1)\cdots(m+n-1).
\]
Теорема доказана.
\textit{(Источник: Липский, гл. 1, §1.2, с. 16--17.)}

\textbf{Пример.} При \(m=3\), \(n=2\) и объектах \(a,b\) число упорядоченных размещений равно
\[
[3]^2=3\cdot4=12.
\]
Это все способы распределить \(a\) и \(b\) по трём ящикам с учётом порядка внутри каждого ящика.
\textit{(Источник: Липский, гл. 1, §1.2, с. 17.)}

Полезны следующие простые соотношения:
\[
[m]_n=(m-n+1)[m]_{n-1},
\]
\[
[m]_n=\frac{m!}{(m-n)!},
\]
\[
[m]^n=[m+n-1]_n.
\]
Последнее равенство показывает, что упорядоченные размещения \(n\) объектов по \(m\) ящикам эквивалентны инъективным размещениям \(n\) объектов по \(m+n-1\) позициям.
\textit{(Источник: Липский, гл. 1, §1.2, с. 17.)}

\paragraph{Перестановки: разложение на циклы.}
Напомним, что перестановкой \(n\)-элементного множества \(X\) называется произвольная взаимно однозначная функция \(f:X\to X\). Обычно перестановка записывается таблицей из двух строк, в которой под элементом \(x\) стоит \(f(x)\). Например, для \(X=\{a,b,c,d\}\) перестановка
\[
f=
\begin{pmatrix}
a & b & c & d\\
d & a & c & b
\end{pmatrix}
\]
задаёт \(f(a)=d\), \(f(b)=a\), \(f(c)=c\), \(f(d)=b\).
\textit{(Источник: Липский, гл. 1, §1.3, с. 17--18.)}

Произвольную перестановку \(f\in S_n\) можно представить ориентированным графом с вершинами \(1,\dots,n\), в котором из \(x\) идёт дуга в \(f(x)\). Так как из каждой вершины выходит ровно одна дуга и в каждую входит ровно одна дуга, граф распадается на непересекающиеся циклы. Это даёт \emph{разложение перестановки на циклы}:
\[
f=[a_0^{(1)}\dots a_{n_1-1}^{(1)}]\cdots[a_0^{(k)}\dots a_{n_k-1}^{(k)}].
\]
Цикл длины \(1\) соответствует неподвижной точке, цикл длины \(2\) называется \emph{транспозицией}. Перестановка имеет \emph{тип} \(1^{\lambda_1}2^{\lambda_2}\dots n^{\lambda_n}\), если в её разложении на циклы ровно \(\lambda_i\) циклов длины \(i\).
\textit{(Источник: Липский, гл. 1, §1.3, с. 19--20.)}

\textbf{Пример.} Перестановка
\[
f=
\begin{pmatrix}
1 & 2 & 3 & 4 & 5 & 6 & 7\\
7 & 5 & 1 & 4 & 2 & 3 & 6
\end{pmatrix}
\]
раскладывается в циклы
\[
f=[1\,7\,6\,3][2\,5][4]
\]
и имеет тип \(1^1 2^1 4^1\).
\textit{(Источник: Липский, гл. 1, §1.3, с. 20.)}

\paragraph{Инверсии и знак перестановки.}
Пусть перестановка задана последовательностью \(\langle a_1,\dots,a_n\rangle\). Пара \(\langle a_i,a_j\rangle\), \(i<j\), называется \emph{инверсией}, если \(a_i>a_j\). Число инверсий обозначается \(I(f)\). \emph{Знак} перестановки определяется равенством
\[
\operatorname{sgn}(f)=(-1)^{I(f)}.
\]
Перестановка называется \emph{чётной}, если \(\operatorname{sgn}(f)=1\), и \emph{нечётной}, если \(\operatorname{sgn}(f)=-1\).
\textit{(Источник: Липский, гл. 1, §1.3, с. 20--21.)}

\textbf{Пример.} Перестановка \(\langle 2,4,3,5,1\rangle\) имеет инверсии
\[
\langle 2,1\rangle,\quad \langle 4,3\rangle,\quad \langle 4,1\rangle,\quad \langle 3,1\rangle,\quad \langle 5,1\rangle,
\]
поэтому \(I(f)=5\) и \(\operatorname{sgn}(f)=-1\); перестановка нечётна.
\textit{(Источник: Липский, гл. 1, §1.3, с. 21.)}

Тождественная перестановка \(\langle 1,\dots,n\rangle\) не имеет инверсий и является чётной. Перестановка \(\langle n,n-1,\dots,1\rangle\) имеет
\[
\frac{n(n-1)}2
\]
инверсий, поэтому её чётность зависит от \(n\): она чётна при \(n=4k\) или \(n=4k+1\) и нечётна в остальных случаях.
\textit{(Источник: Липский, гл. 1, §1.3, с. 21.)}

Справедливы следующие важные свойства.

\textbf{Лемма 1.5.} Произвольную перестановку \(f\in S_n\) можно представить в виде суперпозиции \(I(f)\) транспозиций соседних элементов.

\textbf{Лемма 1.6.} Для любых перестановок \(f,g\in S_n\)
\[
\operatorname{sgn}(fg)=\operatorname{sgn}(f)\operatorname{sgn}(g).
\]

\textbf{Лемма 1.7.} Каждая транспозиция нечётна; знак произвольного цикла длины \(k\) равен \((-1)^{k-1}\).

\textbf{Лемма 1.8.} Если перестановка \(f\) имеет тип \(1^{\lambda_1}2^{\lambda_2}\dots n^{\lambda_n}\), то
\[
\operatorname{sgn}(f)=(-1)^{\lambda_2+\lambda_4+\dots}.
\]
Иными словами, знак перестановки определяется только числом циклов чётной длины.
\textit{(Источник: Липский, гл. 1, §1.3, с. 21--23.)}

На лемме 1.8 основан эффективный алгоритм определения знака перестановки за время \(O(n)\). Он последовательно просматривает позиции \(1,\dots,n\), выделяет циклы перестановки и меняет знак переменной \(s\) при обнаружении цикла чётной длины. По окончании работы \(s=\operatorname{sgn}(f)\).
\textit{(Источник: Липский, гл. 1, §1.3, алгоритм 1.9, с. 23--24.)}

\paragraph{Подмножества и биномиальные коэффициенты.}
Число всех подмножеств \(n\)-элементного множества \(X\) равно \(2^n\). Это устанавливается взаимно однозначным соответствием между подмножествами \(Y\subseteq X\) и бинарными последовательностями \(b_1\dots b_n\), где
\[
b_i=
\begin{cases}
1, & x_i\in Y,\\
0, & x_i\notin Y.
\end{cases}
\]
Число таких последовательностей равно \(2^n\), поэтому
\[
|\mathcal P(X)|=2^n.
\]
\textit{(Источник: Липский, гл. 1, §1.5, с. 34.)}

Число \(k\)-элементных подмножеств \(n\)-элементного множества обозначается
\[
\binom{n}{k}
\]
и называется \emph{биномиальным коэффициентом}. Оно входит в формулу бинома Ньютона:
\[
(x+y)^n=\sum_{k=0}^n \binom{n}{k}x^k y^{n-k}.
\]
\textit{(Источник: Липский, гл. 1, §1.6, с. 39--40.)}

\textbf{Теорема 1.14 (формула для биномиального коэффициента).}
\[
\binom{n}{k}=\frac{[n]_k}{k!}=\frac{n!}{k!(n-k)!}.
\]

\emph{Доказательство.} Каждая инъективная последовательность длины \(k\) с элементами из \(n\)-элементного множества определяет \(k\)-элементное подмножество, состоящее из её элементов. Одно и то же подмножество получается ровно из \(k!\) таких последовательностей, соответствующих всевозможным перестановкам его элементов. Поэтому
\[
\binom{n}{k}=\frac{[n]_k}{k!}=\frac{n!}{k!(n-k)!}.
\]
Теорема доказана.
\textit{(Источник: Липский, гл. 1, §1.6, с. 42.)}

Для биномиальных коэффициентов справедливы важнейшие тождества:
\[
\binom{n}{k}=\binom{n}{n-k},
\]
\[
\binom{n}{k}=\binom{n-1}{k}+\binom{n-1}{k-1},
\]
\[
\sum_{k=0}^n \binom{n}{k}=2^n,
\]
\[
\sum_{k=0}^n k\binom{n}{k}=n2^{n-1}.
\]
Первое выражает симметрию выбора \(k\) элементов и \(n-k\) элементов; второе --- правило треугольника Паскаля; третье --- подсчёт всех подмножеств; четвёртое получается подсчётом пар \(\langle x,K\rangle\), где \(x\in K\subseteq X\), двумя способами.
\textit{(Источник: Липский, гл. 1, §1.6, с. 40--41.)}

Справедливо также \emph{тождество Коши}:
\[
\binom{m+n}{k}=\sum_{s=0}^k \binom{m}{s}\binom{n}{k-s}.
\]
Оно получается выбором \(k\) человек из \(m\) мужчин и \(n\) женщин с классификацией по числу выбранных мужчин. В частности, при \(m=n=k\)
\[
\binom{2n}{n}=\sum_{s=0}^n \binom{n}{s}^2.
\]
\textit{(Источник: Липский, гл. 1, §1.6, с. 41--42.)}

Полезно тождество
\[
\binom{n}{k}\binom{k}{m}=\binom{n}{m}\binom{n-m}{k-m},
\]
которое подсчитывает пары \(\langle K,M\rangle\) подмножеств \(K\subseteq M\subseteq X\) с \(|K|=m\), \(|M|=k\) двумя способами.
\textit{(Источник: Липский, гл. 1, §1.6, с. 42.)}

Биномиальные коэффициенты сначала возрастают, а затем убывают:
\[
\binom{n}{0}<\binom{n}{1}<\dots<\binom{n}{\lfloor n/2\rfloor}
=
\binom{n}{\lceil n/2\rceil}
>\dots>\binom{n}{n}
\]
(при \(n>1\)). Это следует из отношения
\[
\frac{\binom{n}{i+1}}{\binom{n}{i}}=\frac{n-i}{i+1}.
\]
\textit{(Источник: Липский, гл. 1, §1.6, с. 43.)}

\paragraph{Сочетания с повторениями.}
Рассмотрим множество с повторениями
\[
(k*x_1,\dots,k*x_n),
\]
в котором каждый из \(n\) различных элементов может встречаться не более \(k\) раз. Спрашивается, сколько существует \(k\)-элементных подмножеств такого множества с повторениями. Это число равно числу неубывающих последовательностей длины \(k\) с элементами из \(\{1,\dots,n\}\).

\textbf{Теорема 1.15 (о сочетаниях с повторениями).} Число \(k\)-элементных подмножеств множества с повторениями
\[
(k*x_1,\dots,k*x_n)
\]
равно
\[
\binom{n+k-1}{k}.
\]

\emph{Доказательство.} Каждой неубывающей последовательности \(\langle a_1,\dots,a_k\rangle\) с элементами из \(\{1,\dots,n\}\) поставим в соответствие последовательность длины \(k+n-1\), разделив \(n\) групп элементов \(1,\dots,n\) с помощью \(n-1\) нулей-разделителей. Например, при \(n=5\), \(k=4\) последовательности \(\langle 2,2,4,5\rangle\) соответствует
\[
\langle 0,2,2,0,0,4,0,5\rangle.
\]
Такая последовательность однозначно определяется расположением \(n-1\) нулей среди \(k+n-1\) позиций. Число способов выбрать позиции для нулей равно
\[
\binom{k+n-1}{n-1}=\binom{n+k-1}{k}.
\]
Теорема доказана.
\textit{(Источник: Липский, гл. 1, §1.6, с. 43--44.)}

Эту же формулу можно записать в виде
\[
\binom{n+k-1}{k}=\frac{n(n+1)\cdots(n+k-1)}{k!}=\frac{[n]^k}{k!}.
\]
Она показывает, что число \(k\)-элементных сочетаний с повторениями из \(n\) элементов равно числу упорядоченных размещений \(k\) объектов по \(n\) ящикам, делённому на \(k!\).
\textit{(Источник: Липский, гл. 1, §1.6, с. 44--45.)}

\textbf{Пример.} При \(n=3\), \(k=4\) число \(4\)-элементных подмножеств множества с повторениями, в котором каждый из трёх элементов может встречаться не более четырёх раз, равно
\[
\binom{3+4-1}{4}=\binom{6}{4}=15.
\]
\textit{(Источник: Липский, гл. 1, §1.6, с. 43--45.)}

\paragraph{Сводка основных формул.}
Для удобства приведём основные формулы раздела:
\[
\begin{aligned}
&\text{все функции }X\to Y: && m^n;\\
&\text{инъективные функции }X\to Y: && [m]_n=m(m-1)\cdots(m-n+1);\\
&\text{перестановки }n\text{-элементного множества}: && n!;\\
&\text{упорядоченные размещения по ящикам}: && [m]^n=m(m+1)\cdots(m+n-1);\\
&\text{все подмножества}: && 2^n;\\
&\text{\(k\)-элементные подмножества}: && \binom{n}{k}=\frac{n!}{k!(n-k)!};\\
&\text{\(k\)-элементные подмножества с повторениями}: && \binom{n+k-1}{k}.
\end{aligned}
\]
\textit{(Обобщение по: Липский, гл. 1, §1.2--1.3, с. 14--24; §1.5--1.6, с. 34--45.)}

\paragraph{Итог пункта.}
В п. 1.5.1 введены базовые комбинаторные объекты: функции и размещения, инъективные размещения, перестановки, упорядоченные размещения по ящикам, подмножества, сочетания и сочетания с повторениями. Для каждого объекта приведены основные формулы, доказательства и примеры. Показано, что число перестановок равно \(n!\), число \(k\)-элементных подмножеств выражается биномиальным коэффициентом \(\binom{n}{k}\), а сочетания с повторениями считаются по формуле \(\binom{n+k-1}{k}\). Введены разложение перестановки на циклы, инверсии и знак перестановки; сформулированы леммы о знаке и алгоритм его вычисления за \(O(n)\). Приведены важнейшие тождества для биномиальных коэффициентов: симметрия, правило Паскаля, суммы \(\sum \binom{n}{k}=2^n\), \(\sum k\binom{n}{k}=n2^{n-1}\), тождество Коши и тождество для пар подмножеств. Эти результаты служат основой для дальнейшего изучения разбиений множеств и чисел, чисел Стирлинга, производящих функций и метода включений и исключений.
\textit{(Обобщение по: Липский, гл. 1, §1.2--1.3, с. 14--24; §1.5--1.6, с. 34--45.)}

\subsubsection{Разбиения множеств и чисел, числа Стирлинга}

\noindent\textit{Основные источники: Липский В. Комбинаторика для программистов. М.: Мир, 1988. Гл. 1, §1.8--1.11, с. 48--64.}

\paragraph{Разбиения множества.}
Под \emph{разбиением} \(n\)-элементного множества \(X\) на \(k\) блоков понимается семейство
\[
\pi=\{B_1,\dots,B_k\},
\]
такое что
\[
B_1\cup\dots\cup B_k=X,\qquad
B_i\cap B_j=\varnothing\quad(1\le i<j\le k),\qquad
B_i\ne\varnothing\quad(1\le i\le k).
\]
Подмножества \(B_1,\dots,B_k\) называются \emph{блоками} разбиения \(\pi\). Множество всех разбиений \(X\) на \(k\) блоков обозначается \(\Pi_k(X)\), а множество всех разбиений — через \(\Pi(X)\). Очевидно,
\[
\Pi(X)=\Pi_1(X)\cup\Pi_2(X)\cup\dots\cup\Pi_n(X),
\]
причём эти классы попарно не пересекаются.
\textit{(Источник: Липский, гл. 1, §1.8, с. 48.)}

С каждым разбиением \(\pi\) связано отношение эквивалентности
\[
E(\pi)=\bigcup_{B\in\pi}(B\times B),
\]
то есть \(x\) и \(y\) эквивалентны тогда и только тогда, когда они лежат в одном блоке разбиения \(\pi\). Обратно, каждому отношению эквивалентности \(E\) на \(X\) соответствует разбиение
\[
X/E=\{x/E\mid x\in X\},
\]
где \(x/E=\{y\in X\mid xEy\}\). Формулы
\[
E(\pi)=\bigcup_{B\in\pi}(B\times B),
\qquad
X/E=\{x/E\mid x\in X\}
\]
задают взаимно однозначное соответствие между разбиениями множества \(X\) и отношениями эквивалентности на \(X\).
\textit{(Источник: Липский, гл. 1, §1.8, с. 48.)}

Если \(\pi,\sigma\in\Pi(X)\) и каждый блок \(B\in\sigma\) является объединением некоторого числа блоков разбиения \(\pi\), то говорят, что \(\pi\) есть \emph{измельчение} \(\sigma\), и пишут \(\pi\le\sigma\). Например,
\[
(\{1\},\{2,5\},\{4,6\},\{3\})\le(\{1,2,3,5\},\{4,6\}).
\]
Отношение \(\le\) является частичным порядком на \(\Pi(X)\). Более того, \(\Pi(X)\) относительно измельчения образует решётку: для любых \(\pi,\sigma\)
\[
\pi\wedge\sigma=
\{A\cap B\mid A\in\pi,\ B\in\sigma,\ A\cap B\ne\varnothing\},
\]
а \(\pi\vee\sigma\) определяется как наименьшее разбиение, блоки которого являются объединениями блоков из \(\pi\) и \(\sigma\).
\textit{(Источник: Липский, гл. 1, §1.8, с. 48--50.)}

\paragraph{Числа Стирлинга второго рода.}
Число разбиений \(n\)-элементного множества на \(k\) блоков обозначается \(S(n,k)\) и называется \emph{числом Стирлинга второго рода}:
\[
S(n,k)=|\Pi_k(X)|,\qquad |X|=n.
\]
Например, \(S(4,2)=7\), поскольку существуют ровно семь разбиений множества \(\{1,2,3,4\}\) на два блока:
\[
\begin{aligned}
&\{\{1,2,3\},\{4\}\},\quad
\{\{1,2,4\},\{3\}\},\quad
\{\{1,3,4\},\{2\}\},\\
&\{\{1,2\},\{3,4\}\},\quad
\{\{1,3\},\{2,4\}\},\quad
\{\{1,4\},\{2,3\}\},\\
&\{\{1\},\{2,3,4\}\}.
\end{aligned}
\]
Очевидно, \(S(n,k)=0\) при \(k>n\). Принимают также \(S(0,0)=1\): пустое семейство блоков является разбиением пустого множества.
\textit{(Источник: Липский, гл. 1, §1.9, с. 50--51.)}

Для чисел Стирлинга второго рода справедливы рекуррентные соотношения
\[
S(n,k)=S(n-1,k-1)+kS(n-1,k)\quad(0<k<n),
\]
\[
S(n,n)=1\quad(n\ge0),
\]
\[
S(n,0)=0\quad(n>0).
\]
\emph{Доказательство первого соотношения.} Рассмотрим все разбиения множества \(\{1,\dots,n\}\) на \(k\) блоков. Они распадаются на два класса: разбиения, содержащие одноэлементный блок \(\{n\}\), и разбиения, в которых \(n\) входит в блок мощности не меньше двух. Число первых равно \(S(n-1,k-1)\): удалив блок \(\{n\}\), получаем разбиение множества \(\{1,\dots,n-1\}\) на \(k-1\) блоков. Число вторых равно \(kS(n-1,k)\): каждому разбиению множества \(\{1,\dots,n-1\}\) на \(k\) блоков соответствует ровно \(k\) разбиений, полученных добавлением \(n\) в один из \(k\) блоков. Складывая, получаем требуемое.
\textit{(Источник: Липский, гл. 1, §1.9, с. 51.)}

Формулы для \(S(n,k)\) удобно вычислять с помощью «треугольника Стирлинга»: каждое значение, кроме крайних единиц, равно сумме числа, стоящего над ним с левой стороны, и числа, стоящего точно над ним, умноженного на \(k\). Например:
\[
\begin{array}{c|ccccccc}
n\backslash k & 0 & 1 & 2 & 3 & 4 & 5 & 6\\\hline
0 & 1 & 0 & 0 & 0 & 0 & 0 & 0\\
1 & 0 & 1 & 0 & 0 & 0 & 0 & 0\\
2 & 0 & 1 & 1 & 0 & 0 & 0 & 0\\
3 & 0 & 1 & 3 & 1 & 0 & 0 & 0\\
4 & 0 & 1 & 7 & 6 & 1 & 0 & 0\\
5 & 0 & 1 & 15 & 25 & 10 & 1 & 0\\
6 & 0 & 1 & 31 & 90 & 65 & 15 & 1
\end{array}
\]
\textit{(Источник: Липский, гл. 1, §1.9, табл. 1.2, с. 51.)}

Справедливо также соотношение
\[
S(n,k)=\sum_{i=k-1}^{n-1}\binom{n-1}{i}S(i,k-1)\quad(k\ge2).
\]
Оно получается классификацией разбиений по блоку \(B\), содержащему элемент \(n\). Если \(|B|=b\), то блок \(B\) можно выбрать \(\binom{n-1}{b-1}\) способами, а оставшееся множество разбить на \(k-1\) блоков \(S(n-b,k-1)\) способами. Полагая \(i=n-b\), получаем формулу.
\textit{(Источник: Липский, гл. 1, §1.9, с. 52.)}

\paragraph{Числа Белла.}
Число всех разбиений \(n\)-элементного множества называется \emph{числом Белла}:
\[
B_n=|\Pi(X)|=\sum_{k=0}^{n}S(n,k).
\]
Для чисел Белла справедлива рекуррентная зависимость
\[
B_{n+1}=\sum_{i=0}^{n}\binom{n}{i}B_i
\]
(принимают \(B_0=1\)). Она получается классификацией разбиений множества \(\{1,\dots,n+1\}\) по блоку, содержащему элемент \(n+1\): если из этого блока удалить \(n+1\), остаётся разбиение некоторого \(i\)-элементного подмножества, а само подмножество можно выбрать \(\binom{n}{i}\) способами.
\textit{(Источник: Липский, гл. 1, §1.9, с. 52--53.)}

Первые числа Белла:
\[
B_0=1,\quad B_1=1,\quad B_2=2,\quad B_3=5,\quad B_4=15,\quad B_5=52,\quad B_6=203.
\]
\textit{(Источник: Липский, гл. 1, §1.9, табл. 1.3, с. 53.)}

\paragraph{Связь с сюръективными функциями.}
Число функций \(f:X\to Y\), отображающих \(n\)-элементное множество \(X\) на \(k\)-элементное множество \(Y\), обозначим \(s_{n,k}\). Каждой такой функции соответствует её ядро
\[
N(f)=\{f^{-1}(y)\mid y\in Y\},
\]
являющееся разбиением \(X\) на \(k\) блоков. Обратно, каждому разбиению \(\pi\in\Pi_k(X)\) соответствует ровно \(k!\) функций из \(X\) на \(Y\), для которых \(N(f)=\pi\): достаточно взаимно однозначно сопоставить блоки разбиения \(\pi\) элементам множества \(Y\). Поэтому
\[
s_{n,k}=k!\,S(n,k).
\]
Отсюда, в частности, следует явная формула
\[
S(n,k)=\frac{1}{k!}\sum_{i=0}^{k-1}(-1)^i\binom{k}{i}(k-i)^n,
\]
которая будет получена в п. 1.5.4 с помощью принципа включений и исключений.
\textit{(Источник: Липский, гл. 1, §1.9, с. 53; §1.13, с. 73--74.)}

Числа Стирлинга второго рода являются коэффициентами перехода от базиса \(1,x,x^2,\dots\) к базису \(1,[x]_1,[x]_2,\dots\), где
\[
[x]_k=x(x-1)\cdots(x-k+1).
\]
Справедливо тождество
\[
x^n=\sum_{k=0}^{n}S(n,k)[x]_k.
\]
Оно доказывается подсчётом числа всех функций \(f:A\to B\), где \(|A|=n\), \(|B|=x\), двумя способами: слева это \(x^n\), а справа — классификация функций по множеству \(f(A)\) даёт \(\sum_k \binom{x}{k}k!S(n,k)=\sum_k [x]_kS(n,k)\).
\textit{(Источник: Липский, гл. 1, §1.9, теорема 1.18, с. 54.)}

\paragraph{Числа Стирлинга первого рода.}
Числа Стирлинга первого рода \(s(n,k)\) определяются как коэффициенты в разложении
\[
[x]_n=\sum_{k=0}^{n}s(n,k)x^k,
\]
где \([x]_n=x(x-1)\cdots(x-n+1)\). Иными словами, \(s(n,k)\) позволяют перейти от базиса \(1,[x]_1,[x]_2,\dots\) обратно к базису \(1,x,x^2,\dots\). Очевидно, \(s(n,k)=0\) при \(k>n\).
\textit{(Источник: Липский, гл. 1, §1.9, с. 54.)}

Для чисел \(s(n,k)\) справедливы рекуррентные соотношения
\[
s(n,k)=s(n-1,k-1)-(n-1)s(n-1,k),
\]
\[
s(n,n)=1,
\]
\[
s(n,0)=0\quad(n>0).
\]
Первое соотношение получается сравнением коэффициентов при \(x^k\) в равенстве
\[
[x]_n=[x]_{n-1}(x-n+1).
\]
Действительно,
\[
\sum_{k=0}^{n}s(n,k)x^k
=
(x-n+1)\sum_{k=0}^{n-1}s(n-1,k)x^k,
\]
откуда коэффициент при \(x^k\) равен \(s(n-1,k-1)-(n-1)s(n-1,k)\).
\textit{(Источник: Липский, гл. 1, §1.9, с. 54--55.)}

Первые значения \(s(n,k)\):
\[
\begin{array}{c|ccccccc}
n\backslash k & 0 & 1 & 2 & 3 & 4 & 5 & 6\\\hline
0 & 1 & 0 & 0 & 0 & 0 & 0 & 0\\
1 & 0 & 1 & 0 & 0 & 0 & 0 & 0\\
2 & 0 & -1 & 1 & 0 & 0 & 0 & 0\\
3 & 0 & 2 & -3 & 1 & 0 & 0 & 0\\
4 & 0 & -6 & 11 & -6 & 1 & 0 & 0\\
5 & 0 & 24 & -50 & 35 & -10 & 1 & 0\\
6 & 0 & -120 & 274 & -225 & 85 & -15 & 1
\end{array}
\]
\textit{(Источник: Липский, гл. 1, §1.9, табл. 1.4, с. 55.)}

Модуль \(|s(n,k)|\) имеет комбинаторный смысл: это число перестановок \(n\)-элементного множества, имеющих в разложении на циклы ровно \(k\) циклов. Кроме того, справедливо тождество
\[
[x]^n=\sum_{k=0}^{n}|s(n,k)|x^k,
\]
где
\[
[x]^n=x(x+1)\cdots(x+n-1)
\]
--- возрастающий факториал. Матрицы чисел Стирлинга первого и второго рода взаимно обратны:
\[
[s(n,k)]_{0\le n,k\le m}^{-1}=[S(n,k)]_{0\le n,k\le m}.
\]
\textit{(Источник: Липский, гл. 1, §1.9, с. 54--55; задачи 1.34--1.36, с. 81.)}

\paragraph{Разбиения чисел.}
\emph{Разбиением числа} \(n\) на \(k\) слагаемых называется представление
\[
n=a_1+\dots+a_k,
\]
где
\[
a_1\ge a_2\ge\dots\ge a_k>0.
\]
Число таких разбиений обозначается \(P(n,k)\), а число всех разбиений числа \(n\) — через \(P(n)\):
\[
P(n)=\sum_{k=1}^{n}P(n,k)\quad(n>0).
\]
Принимают \(P(0,0)=P(0)=1\). Например, \(P(7)=15\); все разбиения числа \(7\) имеют вид
\[
\begin{aligned}
&7,\quad 6+1,\quad 5+2,\quad 5+1+1,\quad 4+3,\quad 4+2+1,\\
&4+1+1+1,\quad 3+3+1,\quad 3+2+2,\quad 3+2+1+1,\\
&3+1+1+1+1,\quad 2+2+2+1,\quad 2+2+1+1+1,\\
&2+1+1+1+1+1,\quad 1+1+1+1+1+1+1.
\end{aligned}
\]
\textit{(Источник: Липский, гл. 1, §1.11, с. 60--61.)}

Для разбиений чисел часто используется \emph{диаграмма Феррерса}: разбиению \(n=a_1+\dots+a_k\) сопоставляется таблица из \(k\) строк, в \(i\)-й строке которой стоит \(a_i\) точек. Транспонирование диаграммы (перемена ролей строк и столбцов) даёт \emph{сопряжённое разбиение} числа \(n\). Например, разбиению
\[
16=6+4+4+2
\]
соответствует диаграмма, сопряжённая диаграмме разбиения
\[
16=4+4+3+3+1+1.
\]
\textit{(Источник: Липский, гл. 1, §1.11, с. 60--61.)}

\textbf{Теорема 1.20.} Число разбиений числа \(n\) на \(k\) слагаемых равно числу разбиений числа \(n\) с наибольшим слагаемым, равным \(k\).
\[
P(n,k)=P(n,\text{наибольшее слагаемое }=k).
\]
\emph{Доказательство.} Транспонирование диаграммы Феррерса устанавливает взаимно однозначное соответствие между разбиениями числа \(n\) на \(k\) слагаемых и разбиениями числа \(n\) с наибольшим слагаемым \(k\).
\textit{(Источник: Липский, гл. 1, §1.11, теорема 1.20, с. 60.)}

\textbf{Теорема 1.21.} Число разбиений числа \(n\) на попарно различные слагаемые равно числу разбиений числа \(n\) на нечётные слагаемые.
\textit{(Источник: Липский, гл. 1, §1.11, теорема 1.21, с. 61.)}

\emph{Доказательство.} Установим взаимно однозначное соответствие между указанными разбиениями. Рассмотрим разбиение числа \(n\) на нечётные слагаемые \(b_1,\dots,b_p\), где слагаемое \(b_i\) встречается \(r_i\) раз. Пусть
\[
r_i=2^{q_1}+2^{q_2}+\dots\qquad(q_1>q_2>\dots)
\]
--- двоичное представление числа \(r_i\). Заменим \(r_i\) слагаемых \(b_i\) на попарно различные слагаемые
\[
b_i2^{q_1},\quad b_i2^{q_2},\quad\dots
\]
Сумма при этом не меняется. Выполнив такую замену для всех \(i\) и упорядочив слагаемые по невозрастанию, получаем разбиение числа \(n\) на попарно различные слагаемые. Обратное преобразование возможно потому, что каждое натуральное число единственным образом представляется в виде произведения нечётного числа на степень двойки. Например,
\[
26=7+5+5+3+3+1+1+1+1
\]
переходит в
\[
26=10+7+6+2+1.
\]
Теорема доказана.
\textit{(Источник: Липский, гл. 1, §1.11, с. 61--62.)}

Производящая функция для последовательности \(P(0),P(1),P(2),\dots\) имеет вид
\[
\sum_{n=0}^{\infty}P(n)x^n=\prod_{h=1}^{\infty}(1-x^h)^{-1}.
\]
Этот факт будет подробно рассмотрен в п. 1.5.3; здесь отметим лишь, что он даёт удобный аналитический инструмент для изучения разбиений чисел.
\textit{(Источник: Липский, гл. 1, §1.12, с. 68.)}

\paragraph{Алгоритм генерирования разбиений числа.}
Для практических приложений важен не только подсчёт, но и перечисление разбиений числа. В пособии приведён алгоритм 1.22, который генерирует все разбиения числа \(n\) в порядке, обратном лексикографическому. Разбиение хранится в виде массивов \(S[1],\dots,S[d]\) попарно различных слагаемых и \(R[1],\dots,R[d]\), где \(R[i]\) показывает, сколько раз слагаемое \(S[i]\) входит в разбиение. Алгоритм находит каждое следующее разбиение за число шагов, ограниченное константой, не зависящей от \(n\).
\textit{(Источник: Липский, гл. 1, §1.11, алгоритм 1.22, с. 63.)}

\paragraph{Итог пункта.}
В п. 1.5.2 рассмотрены разбиения множеств и чисел. Введены числа Стирлинга второго рода \(S(n,k)\) как числа разбиений \(n\)-элементного множества на \(k\) блоков, доказаны рекуррентные соотношения и приведены таблицы. Введены числа Белла \(B_n\), доказана рекуррентная формула для них. Показана связь \(S(n,k)\) с сюръективными функциями: \(s_{n,k}=k!S(n,k)\), и приведено тождество \(x^n=\sum_k S(n,k)[x]_k\). Определены числа Стирлинга первого рода \(s(n,k)\) через разложение \([x]_n=\sum_k s(n,k)x^k\), доказаны рекуррентные соотношения, указан комбинаторный смысл \(|s(n,k)|\) как числа перестановок с \(k\) циклами. Наконец, введены разбиения чисел \(P(n,k)\), \(P(n)\), диаграммы Феррерса и сопряжённые разбиения; доказаны теоремы о равенстве числа разбиений на \(k\) слагаемых и разбиений с наибольшим слагаемым \(k\), а также о равенстве числа разбиений на попарно различные и на нечётные слагаемые.
\textit{(Обобщение по: Липский, гл. 1, §1.8--1.11, с. 48--64; §1.12, с. 68.)}

\subsubsection{Метод производящих функций}

\noindent\textit{Основные источники: Липский В. Комбинаторика для программистов. М.: Мир, 1988. Гл. 1, §1.12, с. 64--72; задачи 1.43--1.48, с. 82--83.}

\paragraph{Определение производящей функции.}
В комбинаторных задачах на подсчёт числа объектов при наличии ограничений искомым решением часто является последовательность
\[
a_0,a_1,a_2,\dots,
\]
где \(a_k\) — число искомых объектов «размерности» \(k\). Например, если речь идёт о разбиениях числа, можно положить \(a_k=P(k)\); если о подмножествах \(n\)-элементного множества --- \(a_k=\binom{n}{k}\); если о бинарных деревьях --- \(a_k=c_k\) (числа Каталана) и т. д. В таких случаях удобно сопоставить последовательности формальный ряд
\[
A(x)=\sum_{k=0}^{\infty}a_kx^k,
\]
который называется \emph{производящей функцией} для данной последовательности.
\textit{(Источник: Липский, гл. 1, §1.12, с. 64.)}

Слово «формальный» означает, что ряд трактуется только как удобная запись последовательности коэффициентов; вопрос о сходимости при этом не ставится. Мы никогда не вычисляем значение такого ряда в конкретной точке \(x\); мы выполняем над рядами некоторые алгебраические операции, а затем находим коэффициенты при отдельных степенях \(x\). Такой подход позволяет применять к комбинаторным задачам мощную технику, не заботясь о сходимости.
\textit{(Источник: Липский, гл. 1, §1.12, с. 64.)}

\paragraph{Операции над формальными рядами.}
Для произвольных рядов
\[
A(x)=\sum_{k=0}^{\infty}a_kx^k,
\qquad
B(x)=\sum_{k=0}^{\infty}b_kx^k
\]
определяются следующие операции:
\begin{enumerate}[label=\arabic*)]
\item \emph{сложение}:
\[
A(x)+B(x)=\sum_{k=0}^{\infty}(a_k+b_k)x^k;
\]
\item \emph{умножение на число} \(p\) (действительное или комплексное):
\[
pA(x)=\sum_{k=0}^{\infty}pa_kx^k;
\]
\item \emph{произведение Коши} (коротко: произведение):
\[
A(x)\cdot B(x)=\sum_{k=0}^{\infty}c_kx^k,
\]
где
\[
c_k=\sum_{i=0}^{k}a_ib_{k-i}=a_0b_k+a_1b_{k-1}+\dots+a_kb_0.
\]
\end{enumerate}
Если \(a_k=0\) для всех \(k>n\), то ряд \(A(x)\) отождествляется с многочленом \(a_nx^n+\dots+a_0\).
\textit{(Источник: Липский, гл. 1, §1.12, с. 64--65.)}

\paragraph{Связь с аналитическими функциями.}
Из математического анализа известно, что если ряд \(A(x)=\sum a_kx^k\) сходится в некоторой окрестности нуля, то его сумма является аналитической функцией в этой окрестности и
\[
a_k=\frac{A^{(k)}(0)}{k!},\qquad k=0,1,2,\dots
\]
(ряд \(A(x)\) есть не что иное, как ряд Маклорена функции \(A(x)\)). Более того, если \(A(x)\) и \(B(x)\) аналитичны в окрестности нуля, то операции сложения, умножения на число и произведения Коши совпадают с обычными операциями над функциями. Это позволяет отождествлять формальный ряд с аналитической функцией в случаях, когда ряд сходится, и писать, например,
\[
\sum_{k=0}^{\infty}x^k=(1-x)^{-1},
\qquad
\sum_{k=0}^{\infty}\frac{x^k}{k!}=e^x.
\]
Такое отождествление удобно тем, что позволяет пользоваться аналитическими тождествами (дифференцированием, разложением в ряд, почленным интегрированием) при работе с формальными рядами.
\textit{(Источник: Липский, гл. 1, §1.12, с. 65.)}

\paragraph{Простейшие производящие функции.}
Приведём несколько базовых примеров.
\begin{enumerate}[label=\arabic*)]
\item Для биномиальных коэффициентов при фиксированном \(n\):
\[
\sum_{k=0}^{\infty}\binom{n}{k}x^k=\sum_{k=0}^{n}\binom{n}{k}x^k=(1+x)^n.
\]
Это непосредственно следует из формулы бинома Ньютона.
\item Для последовательности \(1,2,4,8,\dots,2^k,\dots\):
\[
\sum_{k=0}^{\infty}2^kx^k=\sum_{k=0}^{\infty}(2x)^k=(1-2x)^{-1}.
\]
\item Для последовательности \(0,1,2,3,\dots\):
\[
\sum_{k=0}^{\infty}kx^k=x\sum_{k=0}^{\infty}kx^{k-1}=x\frac{d}{dx}\sum_{k=0}^{\infty}x^k=x\frac{d}{dx}(1-x)^{-1}=x(1-x)^{-2}.
\]
Здесь использовано почленное дифференцирование, допустимое в окрестности нуля.
\item Для последовательности \(1,1,1,\dots\):
\[
\sum_{k=0}^{\infty}x^k=(1-x)^{-1}.
\]
\end{enumerate}
\textit{(Источник: Липский, гл. 1, §1.12, с. 65--66.)}

\paragraph{Комбинаторная интерпретация произведения.}
Формула \((1+x)^n\) имеет простую комбинаторную интерпретацию. Каждый множитель \((1+x)\) можно трактовать как соответствующий некоторому элементу \(a_i\) множества \(X=\{a_1,\dots,a_n\}\) и как представляющий два возможных числа появлений этого элемента в подмножестве: нуль раз (слагаемое \(x^0=1\)) и один раз (слагаемое \(x^1=x\)). Каждое подмножество множества \(X\) однозначно определяется указанием числа появлений в нём каждого из элементов, то есть выбором одного из слагаемых из каждого множителя произведения \((1+x)\cdots(1+x)\). Определённое таким образом слагаемое разложения даёт вклад в коэффициент при \(x^k\), где \(k\) --- мощность подмножества, равный единице. Это рассуждение естественным образом переносится на множества с повторениями.
\textit{(Источник: Липский, гл. 1, §1.12, с. 66.)}

\paragraph{Множества с повторениями и производящие функции.}
Предположим, что \(X=(2*a_1,3*a_2,1*a_3,4*a_4)\), и обозначим через \(c_k\) число \(k\)-элементных подмножеств этого множества с повторениями. В силу рассуждения, аналогичного приведённому выше, производящая функция для последовательности \(c_0,c_1,c_2,\dots\) равна
\[
\begin{aligned}
\sum_{k=0}^{\infty}c_kx^k
&=(1+x+x^2)(1+x+x^2+x^3)(1+x)(1+x+x^2+x^3+x^4)\\
&=1+4x+9x^2+15x^3+20x^4+22x^5+20x^6\\
&\qquad+15x^7+9x^8+4x^9+x^{10}.
\end{aligned}
\]
Соответствующим подбором \(i\)-го множителя можно накладывать произвольные ограничения на число вхождений элемента \(a_i\). Например, множитель имеет вид \(1+x^3+x^7\), если \(i\)-й элемент может появляться \(0\), \(3\) или \(7\) раз; множитель \(1+x^2+x^4+\dots=(1-x^2)^{-1}\), если этот элемент может появляться произвольное чётное число раз, и т. д. В частности, если не накладывать никаких ограничений на число вхождений элементов \(a_i\), \(1\le i\le n\), то производящая функция имеет вид
\[
(1+x+x^2+\dots)\cdots(1+x+x^2+\dots)=(1-x)^{-n}.
\]
\textit{(Источник: Липский, гл. 1, §1.12, с. 66--67.)}

\paragraph{Доказательство теоремы о сочетаниях с повторениями.}
Напомним (см. п. 1.5.1), что число \(k\)-элементных подмножеств множества с повторениями \((k*x_1,\dots,k*x_n)\) равно \(\binom{n+k-1}{k}\). Покажем, как этот результат получается с помощью производящих функций.

Имеем
\[
(1-x)^{-n}=\sum_{k=0}^{\infty}\binom{n+k-1}{k}x^k.
\]
Действительно, \(k\)-я производная функции \((1-x)^{-n}\) равна
\[
\frac{d^k}{dx^k}(1-x)^{-n}=(-n)(-n-1)\cdots(-n-k+1)(1-x)^{-n-k}(-1)^k=[n]^k(1-x)^{-n-k},
\]
где \([n]^k=n(n+1)\cdots(n+k-1)\). Разлагая \((1-x)^{-n}\) в ряд Маклорена, получаем
\[
(1-x)^{-n}=\sum_{k=0}^{\infty}\frac{[n]^k}{k!}x^k=\sum_{k=0}^{\infty}\binom{n+k-1}{k}x^k.
\]
Это даёт ещё одно независимое доказательство теоремы 1.15.
\textit{(Источник: Липский, гл. 1, §1.12, с. 67.)}

\paragraph{Доказательство тождества Коши.}
Производящие функции удобны и для доказательства тождеств, связанных с биномиальными коэффициентами. Например, тождество Коши
\[
\binom{m+n}{k}=\sum_{s=0}^{k}\binom{m}{s}\binom{n}{k-s}
\]
получается сравнением коэффициентов в обеих частях равенства
\[
\sum_{k=0}^{m+n}\binom{m+n}{k}x^k=(1+x)^{m+n}=(1+x)^m(1+x)^n
=\sum_{i=0}^{m}\binom{m}{i}x^i\cdot\sum_{j=0}^{n}\binom{n}{j}x^j
=\sum_{k=0}^{m+n}\sum_{s=0}^{k}\binom{m}{s}\binom{n}{k-s}x^k.
\]
Здесь использована формула произведения Коши двух многочленов.
\textit{(Источник: Липский, гл. 1, §1.12, с. 67.)}

\paragraph{Разбиения чисел и производящие функции.}
Каждое разбиение \(n=a_1+\dots+a_k\) можно однозначно представить последовательностью \(\langle\lambda_1,\dots,\lambda_n\rangle\), где \(\lambda_i\) обозначает число слагаемых, равных \(i\). Имеем
\[
\sum_{i=1}^{n}i\lambda_i=n,
\]
и каждая последовательность \(\langle\lambda_1,\dots,\lambda_n\rangle\), \(\lambda_i\ge0\), удовлетворяющая этому условию, однозначно определяет разбиение числа \(n\), содержащее \(\lambda_i\) слагаемых, равных \(i\).

Обозначим через \(P_h(n)\) число разбиений числа \(n\) на слагаемые, не превышающие \(h\).

\textbf{Теорема 1.23.} Производящая функция для последовательности \(P_h(0),P_h(1),P_h(2),\dots\) равна
\[
(1+x+x^2+\dots)(1+x^2+x^4+\dots)\cdots(1+x^h+x^{2h}+\dots)
=(1-x)^{-1}(1-x^2)^{-1}\cdots(1-x^h)^{-1}.
\]
\emph{Доказательство.} Согласно определению произведения рядов правая часть является суммой слагаемых вида
\[
x^{\lambda_1}x^{2\lambda_2}\cdots x^{h\lambda_h},
\]
где \(\lambda_i\) --- номер члена, выбранного из \(i\)-го ряда. Поэтому коэффициент при \(x^n\) равен числу последовательностей \(\langle\lambda_1,\dots,\lambda_h\rangle\), таких что
\[
\sum_{i=1}^{h}i\lambda_i=n.
\]
Но это в точности число разбиений числа \(n\) на слагаемые, не превышающие \(h\), то есть \(P_h(n)\).
\textit{(Источник: Липский, гл. 1, §1.12, теорема 1.23, с. 67--68.)}

\textbf{Теорема 1.24.} Производящая функция для последовательности \(P(0),P(1),P(2),\dots\) равна
\[
\prod_{h=1}^{\infty}(1-x^h)^{-1}.
\]
\emph{Доказательство.} Аналогично предыдущему, но число множителей неограниченно. Формально бесконечное произведение определяется как предел частичных произведений; для коэффициента при \(x^n\) достаточно взять первые \(n\) множителей.
\textit{(Источник: Липский, гл. 1, §1.12, теорема 1.24, с. 68.)}

\paragraph{Доказательство теоремы 1.21 через производящие функции.}
Напомним теорему 1.21: число разбиений числа \(n\) на попарно различные слагаемые равно числу разбиений числа \(n\) на нечётные слагаемые. Производящая функция для разбиений на попарно различные слагаемые равна
\[
R(x)=(1+x)(1+x^2)(1+x^3)\cdots(1+x^k)\cdots,
\]
а производящая функция для разбиений на нечётные слагаемые равна
\[
N(x)=(1-x)^{-1}(1-x^3)^{-1}\cdots(1-x^{2k-1})^{-1}\cdots.
\]
Пользуясь зависимостью
\[
1+x^k=\frac{1-x^{2k}}{1-x^k},
\]
получаем
\[
\begin{aligned}
R(x)&=\frac{1-x^2}{1-x}\cdot\frac{1-x^4}{1-x^2}\cdot\frac{1-x^6}{1-x^3}\cdot\dots\\
&=\frac{1}{1-x}\cdot\frac{1}{1-x^3}\cdot\frac{1}{1-x^5}\cdot\dots=N(x).
\end{aligned}
\]
Совпадение производящих функций означает совпадение коэффициентов, то есть равенство чисел разбиений.
\textit{(Источник: Липский, гл. 1, §1.12, с. 68--69.)}

\paragraph{Числа Фибоначчи.}
Числа Фибоначчи \(F_n\) определяются рекуррентным уравнением
\[
F_{n+1}=F_n+F_{n-1},\qquad n\ge1,
\]
с начальными условиями
\[
F_0=F_1=1.
\]
Производящая функция \(F(x)=\sum_{k=0}^{\infty}F_kx^k\) удовлетворяет уравнению
\[
\begin{aligned}
F(x)&=\sum_{k=0}^{\infty}F_kx^k=1+x+\sum_{k=2}^{\infty}(F_{k-2}+F_{k-1})x^k\\
&=1+x+x^2\sum_{k=2}^{\infty}F_{k-2}x^{k-2}+x\sum_{k=2}^{\infty}F_{k-1}x^{k-1}\\
&=1+x+x^2F(x)+x(F(x)-1)=1+(x+x^2)F(x).
\end{aligned}
\]
Отсюда
\[
F(x)=\frac{1}{1-x-x^2}.
\]
\textit{(Источник: Липский, гл. 1, §1.12, с. 69.)}

Найдя корни уравнения \(1-x-x^2=0\), получаем разложение
\[
1-x-x^2=(1-ax)(1-bx),
\]
где
\[
a=\frac{1+\sqrt5}{2},\qquad b=\frac{1-\sqrt5}{2}.
\]
Методом неопределённых коэффициентов находим
\[
\frac{1}{(1-ax)(1-bx)}=\frac{A}{1-ax}+\frac{B}{1-bx}
=\frac{A+B-(Ab+Ba)x}{(1-ax)(1-bx)}.
\]
Сравнивая коэффициенты в числителе, получаем
\[
A+B=1,\qquad Ab+(1-A)a=0,
\]
откуда
\[
A=\frac{a}{a-b},\qquad B=-\frac{b}{a-b}.
\]
Следовательно,
\[
\begin{aligned}
F(x)&=A(1-ax)^{-1}+B(1-bx)^{-1}\\
&=A\sum_{k=0}^{\infty}a^kx^k+B\sum_{k=0}^{\infty}b^kx^k
=\sum_{k=0}^{\infty}\frac{a^{k+1}-b^{k+1}}{a-b}x^k,
\end{aligned}
\]
и окончательно
\[
F_k=\frac{1}{\sqrt5}\left[\left(\frac{1+\sqrt5}{2}\right)^{k+1}-\left(\frac{1-\sqrt5}{2}\right)^{k+1}\right].
\]
Это \emph{формула Бине} для чисел Фибоначчи. Она показывает, что \(F_k\) растёт как \(a^{k+1}/\sqrt5\), где \(a=(1+\sqrt5)/2\) --- золотое сечение.
\textit{(Источник: Липский, гл. 1, §1.12, с. 69--70.)}

\paragraph{Числа Каталана.}
Под \emph{бинарным деревом} с \(n\) вершинами понимается пустое дерево \(T=\varnothing\) при \(n=0\) или тройка
\[
T=\langle L,r,P\rangle,
\]
где \(r\) --- корень дерева, \(L\) --- левое поддерево с \(l\) вершинами, \(P\) --- правое поддерево с \(p\) вершинами, причём \(l+p+1=n\). Два бинарных дерева изоморфны, если они одновременно пусты или их левые и правые поддеревья изоморфны. Обозначим через \(c_k\) число неизоморфных бинарных деревьев с \(k\) вершинами.
\textit{(Источник: Липский, гл. 1, §1.12, с. 70.)}

Из рекурсивного определения следует, что \(c_0=1\) и
\[
c_k=\sum_{s=0}^{k-1}c_sc_{k-1-s},\qquad k>0.
\]
Действительно, дерево с \(k\) вершинами однозначно задаётся числом \(s\) вершин в левом поддереве; число таких деревьев равно \(c_sc_{k-1-s}\), а \(s\) может принимать значения \(0,1,\dots,k-1\).
\textit{(Источник: Липский, гл. 1, §1.12, с. 70.)}

Рассмотрим производящую функцию
\[
C(x)=\sum_{k=0}^{\infty}c_kx^k.
\]
Равенство для \(c_k\) напоминает формулу для коэффициентов произведения Коши:
\[
C(x)C(x)=C^2(x)=\sum_{k=0}^{\infty}\left(\sum_{s=0}^{k}c_sc_{k-s}\right)x^k.
\]
Следовательно,
\[
C(x)=1+xC^2(x),
\]
или
\[
xC^2(x)-C(x)+1=0.
\]
\textit{(Источник: Липский, гл. 1, §1.12, с. 70--71.)}

Рассматривая это равенство как квадратное уравнение относительно \(C(x)\), получаем
\[
C(x)=\frac{1\pm\sqrt{1-4x}}{2x}.
\]
Разложим \(\sqrt{1-4x}=(1-4x)^{1/2}\) в ряд Маклорена. Для \(k>0\)
\[
\begin{aligned}
\frac{d^k}{dx^k}(1-4x)^{1/2}
&=\frac12\left(\frac12-1\right)\left(\frac12-2\right)\cdots\left(\frac12-k+1\right)(1-4x)^{1/2-k}(-4)^k\\
&=2^k(-1)\cdot1\cdot3\cdot5\cdots(2k-3)(1-4x)^{1/2-k}.
\end{aligned}
\]
Следовательно,
\[
\sqrt{1-4x}=1-2\sum_{k=1}^{\infty}\frac{1}{k}\binom{2k-2}{k-1}x^k.
\]
Для получения решения с положительными коэффициентами следует выбрать знак минус:
\[
C(x)=\frac{1-\sqrt{1-4x}}{2x}
=\sum_{k=0}^{\infty}\frac{1}{k+1}\binom{2k}{k}x^k.
\]
Таким образом,
\[
c_k=\frac{1}{k+1}\binom{2k}{k}.
\]
Числа \(c_k\) называются \emph{числами Каталана}. Они возникают во многих комбинаторных задачах: число правильных расстановок скобок в произведении \(n+1\) сомножителей, число способов разбиения выпуклого \((n+2)\)-угольника на треугольники непересекающимися диагоналями и др.
\textit{(Источник: Липский, гл. 1, §1.12, с. 71--72; задачи 1.47--1.48, с. 83.)}

\paragraph{Общий метод решения линейных рекуррентных уравнений.}
Метод производящих функций применим к произвольным линейным рекуррентным уравнениям с постоянными коэффициентами:
\[
C_0a_n+C_1a_{n-1}+\dots+C_ra_{n-r}=0
\]
с начальными условиями, задающими \(a_0,a_1,\dots,a_{r-1}\). Схема решения такова:
\begin{enumerate}[label=\arabic*)]
\item записать производящую функцию \(A(x)=\sum a_nx^n\);
\item умножить рекуррентное соотношение на \(x^n\) и просуммировать по \(n\ge r\);
\item выразить \(A(x)\) через начальные условия, получив рациональную функцию;
\item разложить рациональную функцию на простейшие дроби и разложить каждую дробь в степенной ряд;
\item коэффициент при \(x^n\) даёт явную формулу для \(a_n\).
\end{enumerate}
Именно так были получены формулы Бине для чисел Фибоначчи и формула для чисел Каталана.
\textit{(Источник: Липский, гл. 1, §1.12, с. 70; задача 1.46, с. 83.)}

\paragraph{Экспоненциальные производящие функции.}
Наряду с обычными производящими функциями в комбинаторике используются \emph{экспоненциальные производящие функции}
\[
C(x)=\sum_{k=0}^{\infty}c_k\frac{x^k}{k!}.
\]
Они особенно удобны для последовательностей, в которых важен порядок элементов (например, для перестановок и размещений). Основное свойство экспоненциальных производящих функций состоит в том, что произведению рядов соответствует комбинаторная операция «размеченного произведения»: если \(c_k\) --- число способов составить объект из \(k\) помеченных элементов, а \(d_k\) --- аналогичная величина для другого типа объектов, то коэффициент при \(x^k/k!\) в произведении \(C(x)D(x)\) равен \(\sum_{s=0}^{k}\binom{k}{s}c_sd_{k-s}\), что соответствует выбору \(s\) элементов для объекта первого типа и \(k-s\) элементов для объекта второго типа.
\textit{(Источник: Липский, гл. 1, §1.12, задача 1.43, с. 82.)}

Пусть \(c_k\) --- число различных последовательностей длины \(k\) с элементами из множества \(Y=\{y_1,\dots,y_n\}\), в которых число появлений элемента \(y_i\) принадлежит множеству \(\{m_{i1},m_{i2},\dots\}\). Тогда экспоненциальная производящая функция для последовательности \(c_0,c_1,c_2,\dots\) равна
\[
C(x)=\left(\frac{x^{m_{11}}}{m_{11}!}+\frac{x^{m_{12}}}{m_{12}!}+\dots\right)\cdots\left(\frac{x^{m_{n1}}}{m_{n1}!}+\frac{x^{m_{n2}}}{m_{n2}!}+\dots\right).
\]
Например, если каждый элемент может появляться произвольное число раз, то
\[
C(x)=\left(\sum_{m=0}^{\infty}\frac{x^m}{m!}\right)^n=e^{nx},
\]
откуда \(c_k=n^k\), что совпадает с теоремой 1.1.
\textit{(Источник: Липский, гл. 1, §1.12, задача 1.43, с. 82.)}

\paragraph{Пример: число сюръективных функций.}
Число функций \(f:X\to Y\), отображающих \(n\)-элементное множество на \(m\)-элементное, равно \(s_{n,m}=m!S(n,m)\). С помощью экспоненциальных производящих функций можно получить для \(s_{n,m}\) явную формулу. Рассмотрим экспоненциальную производящую функцию по \(n\) при фиксированном \(m\):
\[
\sum_{n=0}^{\infty}s_{n,m}\frac{x^n}{n!}=(e^x-1)^m.
\]
Действительно, каждый из \(m\) элементов множества \(Y\) должен появиться хотя бы один раз, что даёт множитель \(e^x-1\); элементы независимы, поэтому множители перемножаются. Раскладывая \((e^x-1)^m\) по формуле бинома Ньютона, получаем
\[
(e^x-1)^m=\sum_{i=0}^{m}(-1)^i\binom{m}{i}e^{(m-i)x}
=\sum_{n=0}^{\infty}\left(\sum_{i=0}^{m}(-1)^i\binom{m}{i}(m-i)^n\right)\frac{x^n}{n!}.
\]
Сравнивая коэффициенты, находим
\[
s_{n,m}=\sum_{i=0}^{m}(-1)^i\binom{m}{i}(m-i)^n.
\]
Отсюда
\[
S(n,m)=\frac{1}{m!}\sum_{i=0}^{m}(-1)^i\binom{m}{i}(m-i)^n.
\]
Это та же формула, которая будет получена в п. 1.5.4 с помощью принципа включений и исключений.
\textit{(Источник: Липский, гл. 1, §1.12, с. 67; §1.13, с. 73--74.)}

\paragraph{Итог пункта.}
В п. 1.5.3 введено понятие производящей функции как формального ряда, сопоставленного последовательности. Определены операции сложения, умножения на число и произведения Коши; установлена связь с аналитическими функциями. Приведены простейшие производящие функции: \((1+x)^n\), \((1-2x)^{-1}\), \(x(1-x)^{-2}\), \((1-x)^{-1}\). Показано, как производящие функции применяются к множествам с повторениями, к доказательству теоремы о сочетаниях с повторениями и тождества Коши. С помощью производящих функций доказаны теоремы о разбиениях чисел (теоремы 1.23, 1.24), получены формулы Бине для чисел Фибоначчи и явная формула для чисел Каталана. Описан общий метод решения линейных рекуррентных уравнений. Введены экспоненциальные производящие функции и показано их применение к подсчёту сюръективных функций и чисел Стирлинга второго рода.
\textit{(Обобщение по: Липский, гл. 1, §1.12, с. 64--72; задачи 1.43--1.48, с. 82--83.)}

\subsubsection{Метод включений и исключений}

\noindent\textit{Основные источники: Липский В. Комбинаторика для программистов. М.: Мир, 1988. Гл. 1, §1.13, с. 72--76.}

\paragraph{Постановка задачи.}
Пусть дано конечное множество \(X\) и набор подмножеств \(A_1,\dots,A_n\subseteq X\). Требуется найти мощность объединения \(A_1\cup\dots\cup A_n\) или, что эквивалентно, мощность множества элементов, не принадлежащих ни одному из \(A_i\). Метод включений и исключений даёт точную формулу для этой мощности через мощности всевозможных пересечений.

\paragraph{Формула включений и исключений.}
Для любых множеств \(A_1,\dots,A_n\) справедливо равенство
\[
\left|\bigcup_{i=1}^n A_i\right|
=
\sum_{i=1}^n |A_i|
-
\sum_{1\le i<j\le n} |A_i\cap A_j|
+
\sum_{1\le i<j<k\le n} |A_i\cap A_j\cap A_k|
-
\dots
+
(-1)^{n-1} |A_1\cap\dots\cap A_n|.
\]
Эквивалентная формулировка для числа элементов, не обладающих ни одним из свойств \(A_i\):
\[
\left|X \setminus \bigcup_{i=1}^n A_i\right|
=
|X|
-
\sum_{i=1}^n |A_i|
+
\sum_{1\le i<j\le n} |A_i\cap A_j|
-
\dots
+
(-1)^n |A_1\cap\dots\cap A_n|.
\]
Знак перед членом, соответствующим пересечению \(k\) множеств, равен \((-1)^{k-1}\) (в первой формуле) или \((-1)^k\) (во второй).

\paragraph{Доказательство (индукция по \(n\)).}
Докажем первую формулу индукцией по числу множеств \(n\). При \(n=1\) формула тривиальна: \(|A_1|=|A_1|\). Предположим, что формула верна для \(n-1\) множеств. Рассмотрим объединение
\[
A_1\cup\dots\cup A_n = (A_1\cup\dots\cup A_{n-1}) \cup A_n.
\]
По формуле для двух множеств:
\[
|(A_1\cup\dots\cup A_{n-1}) \cup A_n|
=
|A_1\cup\dots\cup A_{n-1}| + |A_n| - |(A_1\cup\dots\cup A_{n-1}) \cap A_n|.
\]
По индукционному предположению
\[
|A_1\cup\dots\cup A_{n-1}|
=
\sum_{i=1}^{n-1}|A_i|
-
\sum_{1\le i<j\le n-1}|A_i\cap A_j|
+
\dots
+
(-1)^{n-2}|A_1\cap\dots\cap A_{n-1}|.
\]
Кроме того,
\[
(A_1\cup\dots\cup A_{n-1}) \cap A_n
=
(A_1\cap A_n) \cup \dots \cup (A_{n-1}\cap A_n).
\]
Применяя к этому объединению \(n-1\) множеств индукционное предположение, получаем
\[
\begin{aligned}
|(A_1\cup\dots\cup A_{n-1}) \cap A_n|
&=
\sum_{i=1}^{n-1}|A_i\cap A_n|
-
\sum_{1\le i<j\le n-1}|A_i\cap A_j\cap A_n|
+
\dots \\
&\qquad
+
(-1)^{n-2}|A_1\cap\dots\cap A_{n-1}\cap A_n|.
\end{aligned}
\]
Подставляя эти выражения в формулу для \(|(A_1\cup\dots\cup A_{n-1})\cup A_n|\) и группируя члены, получаем требуемое равенство. Доказательство завершено.
\textit{(Источник: Липский, гл. 1, §1.13, с. 73.)}

\paragraph{Комбинаторный смысл.}
Формулу можно интерпретировать следующим образом. Если элемент принадлежит ровно \(r\) из множеств \(A_i\) (\(1\le r\le n\)), то в правой части первой формулы он учитывается
\[
\binom{r}{1} - \binom{r}{2} + \binom{r}{3} - \dots + (-1)^{r-1}\binom{r}{r} = 1
\]
раз. Если же элемент не принадлежит ни одному из \(A_i\), он не учитывается. Таким образом, каждый элемент объединения учитывается ровно один раз, а элементы вне объединения — ноль раз. Это даёт альтернативное доказательство.
\textit{(Источник: Липский, гл. 1, §1.13, с. 72.)}

\paragraph{Пример: число сюръективных функций.}
Пусть \(X\) — множество из \(n\) элементов, \(Y\) — множество из \(m\) элементов. Требуется найти число функций \(f:X\to Y\), отображающих \(X\) на \(Y\) (сюръективных). Обозначим это число через \(s_{n,m}\). Всего функций из \(X\) в \(Y\) ровно \(m^n\). Для каждого \(y\in Y\) пусть \(A_y\) — множество функций \(f:X\to Y\), в образе которых не содержится \(y\), то есть \(y\notin f(X)\). Тогда функция не является сюръективной тогда и только тогда, когда она принадлежит хотя бы одному из множеств \(A_y\). По формуле включений и исключений:
\[
s_{n,m} = m^n - \left|\bigcup_{y\in Y} A_y\right|.
\]
Для любого подмножества \(S\subseteq Y\) мощности \(i\) пересечение \(\bigcap_{y\in S} A_y\) состоит из функций, не принимающих значений из \(S\), то есть функций из \(X\) в \(Y\setminus S\). Число таких функций равно \((m-i)^n\). Количество подмножеств \(S\) мощности \(i\) равно \(\binom{m}{i}\). Поэтому
\[
\left|\bigcup_{y\in Y} A_y\right|
=
\sum_{i=1}^{m} (-1)^{i-1} \binom{m}{i} (m-i)^n.
\]
Следовательно,
\[
s_{n,m} = \sum_{i=0}^{m} (-1)^i \binom{m}{i} (m-i)^n.
\]
В частности, при \(m=n\) получаем число перестановок \(n\)-элементного множества, но эта формула даёт также число сюръективных функций для любых \(n,m\).
\textit{(Источник: Липский, гл. 1, §1.13, с. 73--74.)}

\paragraph{Числа Стирлинга второго рода.}
Известно, что \(s_{n,m} = m! S(n,m)\), где \(S(n,m)\) — число Стирлинга второго рода (число разбиений \(n\)-элементного множества на \(m\) блоков). Поэтому
\[
S(n,m) = \frac{1}{m!} \sum_{i=0}^{m} (-1)^i \binom{m}{i} (m-i)^n.
\]
Это даёт явную формулу для чисел Стирлинга второго рода. Например, при \(m=2\):
\[
S(n,2) = \frac{1}{2}\left(2^n - 2\cdot 1^n\right) = 2^{n-1} - 1.
\]
При \(n=4\): \(S(4,2)=7\), что совпадает с ранее найденным значением.
\textit{(Источник: Липский, гл. 1, §1.13, с. 74.)}

\paragraph{Пример: беспорядки (derangements).}
Беспорядком (или инверсией) на множестве \(\{1,\dots,n\}\) называется перестановка \(f\), не имеющая неподвижных точек: \(f(i)\ne i\) для всех \(i\). Обозначим число беспорядков через \(D_n\). Всего перестановок \(n!\). Для каждого \(i\) пусть \(A_i\) — множество перестановок, оставляющих \(i\) на месте: \(f(i)=i\). Тогда беспорядки — это перестановки, не принадлежащие ни одному из \(A_i\). По формуле включений и исключений:
\[
D_n = n! - \sum_{i=1}^n |A_i| + \sum_{1\le i<j\le n} |A_i\cap A_j| - \dots + (-1)^n |A_1\cap\dots\cap A_n|.
\]
Множество \(A_i\) состоит из перестановок, фиксирующих \(i\); таких перестановок \((n-1)!\). Пересечение \(A_i\cap A_j\) состоит из перестановок, фиксирующих \(i\) и \(j\); таких перестановок \((n-2)!\), и так далее. Пересечение \(k\) множеств имеет мощность \((n-k)!\). Количество способов выбрать \(k\) индексов равно \(\binom{n}{k}\). Поэтому
\[
D_n = \sum_{k=0}^{n} (-1)^k \binom{n}{k} (n-k)! = n! \sum_{k=0}^{n} \frac{(-1)^k}{k!}.
\]
Сумма в скобках является частичной суммой ряда \(e^{-1} = \sum_{k=0}^{\infty} (-1)^k/k!\). Поэтому \(D_n \approx n!/e\), и доля беспорядков среди всех перестановок стремится к \(e^{-1} \approx 0.36788\). Например, \(D_4 = 9\), \(D_5 = 44\).
\textit{(Источник: Липский, гл. 1, §1.13, с. 75--76.)}

\paragraph{Пример: подмножества множества с повторениями с ограничениями.}
Рассмотрим задачу о числе \(k\)-элементных подмножеств множества с повторениями \(X = (k_1 * a_1, \dots, k_n * a_n)\), где кратность элемента \(a_i\) не превосходит \(k_i\). Эта задача эквивалентна подсчёту числа целочисленных решений системы
\[
x_1+\dots+x_n = k, \qquad 0\le x_i\le k_i \quad (i=1,\dots,n),
\]
где \(x_i\) — кратность элемента \(a_i\) в подмножестве. Без ограничений \(x_i\le k_i\) число решений равно \(\binom{n+k-1}{k}\) (сочетания с повторениями). Чтобы учесть ограничения, применим метод включений и исключений. Пусть \(A_i\) — множество решений, в которых \(x_i > k_i\). Тогда искомое число равно
\[
N = \binom{n+k-1}{k} - \left|\bigcup_{i=1}^n A_i\right|.
\]
Для подсчёта \(|A_i|\) сделаем замену \(x_i' = x_i - (k_i+1)\ge 0\). Тогда уравнение принимает вид
\[
x_1+\dots+x_{i-1}+x_i'+x_{i+1}+\dots+x_n = k - (k_i+1),
\]
и число неотрицательных решений равно
\[
\binom{n + (k - k_i -1) -1}{k - k_i -1} = \binom{n+k-k_i-2}{k-k_i-1},
\]
если \(k > k_i\), и \(0\) в противном случае. Аналогично для пересечений \(A_i\cap A_j\) делаем замены для двух переменных и получаем \(\binom{n+k-k_i-k_j-3}{k-k_i-k_j-2}\) и т.д. Общая формула:
\[
N = \sum_{S\subseteq\{1,\dots,n\}} (-1)^{|S|} \binom{n+k - \sum_{i\in S}(k_i+1) - 1}{k - \sum_{i\in S}(k_i+1)},
\]
где биномиальный коэффициент считается равным нулю, если верхний индекс меньше нижнего или нижний отрицателен.
\textit{(Источник: Липский, гл. 1, §1.13, с. 74--75.)}

\paragraph{Пример: число простых чисел.}
Хотя этот пример обычно приводится в учебниках по теории чисел, он хорошо иллюстрирует метод. Пусть \(\pi(x)\) — количество простых чисел, не превосходящих \(x\). Используя решето Эратосфена, можно выразить \(\pi(n) - \pi(\sqrt{n})\) через включения-исключения. Если \(p_1,\dots,p_k\) — все простые числа, не превосходящие \(\sqrt{n}\), то число чисел, не превосходящих \(n\) и не делящихся ни на одно из \(p_i\), равно
\[
n - \sum_i \left\lfloor\frac{n}{p_i}\right\rfloor + \sum_{i<j} \left\lfloor\frac{n}{p_ip_j}\right\rfloor - \dots + (-1)^k \left\lfloor\frac{n}{p_1\dots p_k}\right\rfloor.
\]
Это число равно \(\pi(n) - \pi(\sqrt{n}) + 1\) (нужно учесть, что 1 не является простым, а сами простые \(p_i\) не делятся на другие \(p_j\), поэтому они учитываются). В итоге получается формула
\[
\pi(n) = \pi(\sqrt{n}) - 1 + \sum_{S\subseteq\{1,\dots,k\}} (-1)^{|S|} \left\lfloor\frac{n}{\prod_{i\in S}p_i}\right\rfloor.
\]
Эта формула имеет скорее теоретическое значение, так как число слагаемых растёт экспоненциально с \(k\), но она демонстрирует универсальность метода.
\textit{(Источник: Липский, гл. 1, §1.13, задача 1.49, с. 83.)}

\paragraph{Вычислительные аспекты.}
Формула включений и исключений содержит \(2^n\) слагаемых, поэтому при больших \(n\) прямое вычисление становится неэффективным. Однако во многих задачах структура пересечений позволяет сгруппировать слагаемые и получить более компактные выражения (как в примерах с сюръективными функциями и беспорядками). Кроме того, часто достаточно вычислить несколько первых членов, чтобы получить хорошее приближение, поскольку знаки чередуются и вклад последующих членов быстро убывает.

\paragraph{Итог пункта.}
В п. 1.5.4 сформулирован и доказан принцип включений и исключений. Приведены его комбинаторная интерпретация и основные примеры применения: подсчёт числа сюръективных функций, чисел Стирлинга второго рода, беспорядков, подмножеств множества с повторениями с ограничениями, а также числа простых чисел. Показано, что метод даёт явные формулы, удобные для вычислений и анализа, и что он тесно связан с другими комбинаторными понятиями (числами Стирлинга, числами Белла, производящими функциями).
\textit{(Обобщение по: Липский, гл. 1, §1.13, с. 72--76; задачи 1.49--1.50, с. 83.)}

\subsubsection{Краткое резюме}

Комбинаторный анализ изучает способы подсчёта числа объектов, удовлетворяющих заданным ограничениям. Основные понятия — размещения, перестановки, сочетания (с повторениями и без) — выражаются через факториалы и биномиальные коэффициенты. Разбиения множеств приводят к числам Стирлинга и Белла, разбиения чисел — к функции \(P(n)\). Метод производящих функций позволяет сводить комбинаторные задачи к алгебраическим операциям над формальными рядами и получать явные формулы (числа Фибоначчи, Каталана и др.). Метод включений и исключений даёт универсальный способ подсчёта мощности объединения множеств и применяется к задачам о сюръективных функциях, беспорядках и ограниченных композициях.
\textit{(Обобщение по: Липский, гл. 1, §1.1--1.13, с. 9--76.)}

\noindent\textit{Источники к пункту 1.4:}
\begin{itemize}
\item Липский В. Комбинаторика для программистов. М.: Мир, 1988. Гл. 1, §1.1--1.13, с. 9--76.
\end{itemize}

\subsection{Коды с исправлением ошибок. Алфавитное кодирование. Методы сжатия информации}

\noindent\textit{Основной источник: Прокопенко Н.Ю. Теория кодирования: учеб. пособие / Н.Ю. Прокопенко; Нижегор. гос. архитектур.-строит. ун-т. – Н. Новгород: ННГАСУ, 2023. – 94 с.}

\subsubsection{Алфавитное кодирование и взаимная однозначность}

\noindent\textit{Основные источники: Прокопенко Н.Ю. Теория кодирования: учеб. пособие. Н. Новгород: ННГАСУ, 2023. Разд. 2, с. 17--22; разд. 3, с. 23--28.}

\paragraph{Алфавиты, слова и кодирование.}
Пусть задано множество \(A=\{a_1,\dots,a_n\}\), которое называется \emph{алфавитом сообщений}, и множество \(B=\{b_1,\dots,b_m\}\), называемое \emph{кодирующим алфавитом}. Элементы алфавита называются \emph{буквами}. Конечная последовательность букв алфавита \(A\) называется \emph{словом} в алфавите \(A\). Множество всех слов в алфавите \(A\) обозначается \(A^*\), множество непустых слов — \(A^+\). Через \(\lambda\) обозначается пустое слово:
\[
A^*=A^+\cup\{\lambda\}.
\]
\emph{Длиной слова} называется число букв в нём. Если слово \(\alpha\) имеет длину \(k\), пишут \(|\alpha|=k\). Для пустого слова \(|\lambda|=0\). Соединение слов \(\alpha\) и \(\beta\) обозначается \(\alpha\beta\); через \(\alpha^n\) обозначается соединение \(n\) одинаковых слов, причём \(\alpha^0=\lambda\).
\textit{(Источник: Прокопенко, разд. 2, с. 17.)}

\emph{Кодированием} называют отображение
\[
F\colon S\to B^*,
\]
где \(S\subseteq A^*\) — некоторое множество сообщений. Слово \(b=F(a)\) называется \emph{кодом} слова \(a\). Обратное отображение \(F^{-1}\) (если оно существует) называется \emph{декодированием}. Кодирование \(F\) называется \emph{взаимно однозначным} (или \emph{однозначно декодируемым}), если каждый код сообщения является кодом ровно одного сообщения. В этом случае код называется \emph{свободным}.
\textit{(Источник: Прокопенко, разд. 2, с. 17--18.)}

Если все кодовые слова имеют одинаковую длину, код называется \emph{равномерным} (или \emph{блочным}). В противном случае код называется \emph{неравномерным}. Например, код \(e_4(n)\), сопоставляющий числам \(0,1,2,\dots\) их четырёхразрядные двоичные записи, является равномерным; код \(e(n)\), в котором числа записываются двоичными словами без фиксированной длины, равномерным не является.
\textit{(Источник: Прокопенко, разд. 2, с. 17.)}

\paragraph{Алфавитное (побуквенное) кодирование.}
Пусть задано отображение \(\Sigma\) букв алфавита \(A\) в множество \(B^*\):
\[
\Sigma:\quad
\begin{cases}
a_1\to B_1,\\
a_2\to B_2,\\
\dots\\
a_r\to B_r,
\end{cases}
\]
где \(B_i\in B^*\). \emph{Алфавитным кодированием}, задаваемым схемой \(\Sigma\), называется кодирование \(F_\Sigma\), удовлетворяющее двум условиям:
\begin{enumerate}[label=\arabic*)]
\item \(F_\Sigma(a_i)=B_i\), \(i=1,\dots,r\);
\item \(F_\Sigma(a_{i_1}a_{i_2}\dots a_{i_k})=B_{i_1}B_{i_2}\dots B_{i_k}\).
\end{enumerate}
Иными словами, код слова получается последовательным приписыванием кодов его букв. Слова \(B_i\) называются \emph{элементарными кодами}. Множество
\[
C(\Sigma)=\{B_1,\dots,B_r\}
\]
называется \emph{кодом алфавита} в схеме \(\Sigma\).
\textit{(Источник: Прокопенко, разд. 2, с. 19.)}

\emph{Длиной кода} называется число знаков в коде без учёта пробелов. \emph{Спектром кода} называется упорядоченная по невозрастанию последовательность длин кодовых слов:
\[
L(V)=\langle l_1,l_2,\dots,l_r\rangle,\qquad l_1\ge l_2\ge\dots\ge l_r.
\]
Код называется \emph{разделимым}, если любое слово, составленное из элементарных кодов, единственным образом разлагается на элементарные коды. Иными словами, из равенства
\[
B_{i_1}B_{i_2}\dots B_{i_k}=B_{j_1}B_{j_2}\dots B_{j_l}
\]
должно следовать \(k=l\) и \(i_t=j_t\) для всех \(t=1,\dots,k\).
\textit{(Источник: Прокопенко, разд. 2, с. 19.)}

Алфавитное кодирование является взаимно однозначным тогда и только тогда, когда оно задаётся с помощью разделимого кода. Из определения разделимого кода следует, что все слова разделимого кода различны и не пусты.
\textit{(Источник: Прокопенко, разд. 2, с. 19--20.)}

\paragraph{Примеры.}
\begin{enumerate}[label=\arabic*)]
\item Пусть \(A=\{a,b,c,d\}\), \(B=\{0,1\}\), а схема кодирования имеет вид
\[
a\to 01,\qquad b\to 100,\qquad c\to 101,\qquad d\to 0.
\]
Это кодирование не является взаимно однозначным, а значит, код не является свободным. Например, слово \(0100\) можно декодировать как \(db\) или как \(add\); слово \(0010100\) допускает декодирования \(ddcdd\), \(daadd\), \(dadb\).
\item Азбука Морзе не является свободным кодом без дополнительных пауз между буквами и словами. Поэтому при передаче сообщений с помощью азбуки Морзе используют пробелы (паузы), которые неформально отделяют буквы друг от друга.
\end{enumerate}
\textit{(Источник: Прокопенко, разд. 2, с. 18.)}

\paragraph{Неравенство Макмиллана.}
Для того чтобы схема алфавитного кодирования была разделимой в \(q\)-буквенном кодирующем алфавите, необходимо, чтобы длины элементарных кодов удовлетворяли неравенству Макмиллана:
\[
\frac{1}{q^{|B_1|}}+\frac{1}{q^{|B_2|}}+\dots+\frac{1}{q^{|B_r|}}\le 1.
\]
В частности, для двоичного кодирующего алфавита \(B=\{0,1\}\):
\[
\frac{1}{2^{|B_1|}}+\frac{1}{2^{|B_2|}}+\dots+\frac{1}{2^{|B_r|}}\le 1.
\]
\textit{(Источник: Прокопенко, разд. 2, с. 21.)}

Неравенство Макмиллана является \emph{необходимым} условием разделимости, но не является достаточным. Например, схема
\[
a\to 0,\qquad b\to 01
\]
является разделимой, хотя и не является префиксной: \(1/2+1/4=3/4\le 1\). В то же время азбука Морзе не удовлетворяет неравенству Макмиллана, поэтому без дополнительных пауз она не является разделимой.
\textit{(Источник: Прокопенко, разд. 2, с. 21--22.)}

Кроме того, справедлива теорема: если для чисел \(l_1,\dots,l_n\) выполнено неравенство Макмиллана
\[
\frac{1}{q^{l_1}}+\frac{1}{q^{l_2}}+\dots+\frac{1}{q^{l_n}}\le 1,
\]
то существует префиксный код \(V=\{v_1,\dots,v_n\}\) в алфавите \(B=\{b_1,\dots,b_m\}\) со спектром длин \(L(V)=\langle l_1,\dots,l_n\rangle\).
\textit{(Источник: Прокопенко, разд. 4.2, с. 33--34.)}

\paragraph{Проблема взаимной однозначности.}
Одним из основных вопросов теории кодирования является проблема взаимной однозначности: по коду \(\beta\) сообщения \(\alpha\) требуется однозначно восстановить \(\alpha\). Если решать эту задачу исходя непосредственно из определения, то нужно перебрать бесконечное число слов, так как для каждого кода нужно установить, допускает он однозначное декодирование или нет. Для эффективного решения этой задачи используется алгоритм Маркова, основанный на понятии нетривиального разложения элементарных кодов.
\textit{(Источник: Прокопенко, разд. 2, с. 20; разд. 3, с. 23.)}

\paragraph{Нетривиальные разложения.}
Пусть \(C(\Sigma)\) — алфавитный код. Обозначим через \(P_\Sigma\) множество слов \(\beta'\), каждое из которых является \emph{собственным префиксом} некоторого кодового слова \(B_i\). Через \(S_\Sigma\) обозначим множество слов \(\beta''\), каждое из которых является \emph{собственным суффиксом} некоторого кодового слова \(B_i\).

\textbf{Определение.} Представление элементарного кода \(B_i\) в виде
\[
B_i=\beta'\omega\beta'',
\]
где слово \(\beta'\) не может оканчиваться на элементарный код, а слово \(\beta''\) не может начинаться с элементарного кода, называется \emph{нетривиальным разложением} элементарного кода. При этом одно из слов \(\beta'\), \(\omega\), \(\beta''\) может быть пустым. Для наглядности слова в разложении отделяют точками:
\[
B_i=\beta'.\omega.\beta''.
\]
\textit{(Источник: Прокопенко, разд. 3, с. 23.)}

\paragraph{Теорема Маркова.}
Алфавитное кодирование со схемой \(\Sigma\) не обладает свойством взаимной однозначности тогда и только тогда, когда граф \(G_\Sigma\) содержит ориентированный цикл (контур), проходящий через вершину \(\lambda\). Иными словами, алфавитное кодирование является взаимно однозначным (код \(C(\Sigma)\) свободен) тогда и только тогда, когда граф \(G_\Sigma\) не содержит ориентированного цикла, проходящего через вершину \(\lambda\).
\textit{(Источник: Прокопенко, разд. 3, с. 23.)}

Если схема \(\Sigma\) не обладает свойством взаимной однозначности, то слово из \(B^*\), допускающее два декодирования, легко находится по графу \(G_\Sigma\). Для этого нужно рассмотреть ориентированный цикл, проходящий через вершину \(\lambda\), начиная с \(\lambda\), и выписать последовательно все слова, приписанные рёбрам и вершинам, входящим в этот цикл.
\textit{(Источник: Прокопенко, разд. 3, с. 23.)}

\paragraph{Алгоритм построения графа \(G_\Sigma\).}
Пусть \(V=P_\Sigma\cap S_\Sigma\cup\{\lambda\}\) — множество вершин графа. Вершины \(\beta'\) и \(\beta''\) соединяются ориентированным ребром, если для некоторого элементарного кода \(B_i\) существует нетривиальное разложение
\[
B_i=\beta'.\omega.\beta'',
\]
где \(\beta'\in P_\Sigma\), \(\beta''\in S_\Sigma\), \(\omega\in C^*(\Sigma)\). Ребру приписывается значение \(\omega\). Множество всех рёбер обозначается \(R\), множество значений рёбер — \(m(R)\). Полученный взвешенный граф обозначается \(G_\Sigma\).

Алгоритм построения графа:
\begin{enumerate}[label=\arabic*)]
\item Для каждого элементарного кода выписать все нетривиальные разложения \(B_i=\beta'.\omega.\beta''\).
\item Выписать множество \(P_\Sigma\), состоящее из слов \(\beta'\), которые входят в качестве начал в нетривиальные разложения элементарных кодов.
\item Выписать множество \(S_\Sigma\), состоящее из всех слов \(\beta''\), которые являются окончаниями нетривиальных разложений элементарных кодов.
\item Составить множество \(V=P_\Sigma\cap S_\Sigma\cup\{\lambda\}\).
\item Выписать все разложения элементарных кодов, связанных с множеством \(V\), то есть все разложения вида
\[
B_i=\beta'\beta_{i_1}\dots\beta_{i_k}\beta'',
\]
где \(\beta',\beta''\in V\), а \(k\) может быть равно \(0\). Если \(k=0\), то \(\beta_{i_1}\dots\beta_{i_k}=\lambda\).
\item По разложениям, полученным в п. 5, построить ориентированный граф \(G_\Sigma\): вершины графа отождествляются с элементами множества \(V\); пара вершин \(\beta'\) и \(\beta''\) соединяется ориентированным ребром в том и только в том случае, если существует разложение \(B_i=\beta'\beta_{i_1}\dots\beta_{i_k}\beta''\). При этом ребру \((\beta',\beta'')\) приписывается слово \(\beta_{i_1}\dots\beta_{i_k}\).
\item По полученному графу \(G_\Sigma\) проверить, обладает ли исходная схема кодирования свойством взаимной однозначности, используя теорему Маркова.
\end{enumerate}
\textit{(Источник: Прокопенко, разд. 3, с. 23--24.)}

\paragraph{Пример 5.}
Пусть \(A=\{a_1,a_2,a_3\}\), \(B=\{0,1\}\), а код задан схемой
\[
\Sigma:\quad
\begin{cases}
a_1\to 1,\\
a_2\to 010,\\
a_3\to 101.
\end{cases}
\]
Выпишем нетривиальные разложения элементарных кодов:
\[
B_1=1
\]
не имеет нетривиальных разложений;
\[
B_2=010=0.B_1.0=01.\lambda.0=0.\lambda.10;
\]
\[
B_3=101=\lambda.B_1.01=10.B_1.\lambda=1.\lambda.01=10.\lambda.1.
\]
Тогда
\[
P_\Sigma=\{0,01,\lambda,10,1\},\qquad
S_\Sigma=\{0,01,\lambda,10,1\},
\]
\[
V=P_\Sigma\cap S_\Sigma\cup\{\lambda\}=\{0,1,01,\lambda,10\}.
\]
Соответствующий граф \(G_\Sigma\) содержит ориентированный цикл, проходящий через вершину \(\lambda\). Поэтому код не является свободным. По графу можно построить слово, допускающее два декодирования, например слово
\[
1010101.
\]
Оно допускает две расшифровки:
\[
101\,010\,1=a_3a_2a_1,
\qquad
1\,010\,101=a_1a_2a_3.
\]
\textit{(Источник: Прокопенко, разд. 3, с. 25--26.)}

\paragraph{Пример 6.}
Пусть \(A=\{a_1,a_2,a_3,a_4,a_5\}\), \(B=\{0,1\}\), а код задан схемой
\[
\Sigma:\quad
\begin{cases}
a_1\to 1,\\
a_2\to 01,\\
a_3\to 100,\\
a_4\to 0100,\\
a_5\to 0000.
\end{cases}
\]
Выпишем нетривиальные разложения:
\[
B_1=1
\]
не имеет нетривиальных разложений;
\[
B_2=01=0.B_1.\lambda=0.\lambda.1;
\]
\[
B_3=100=\lambda.B_1.00=1.\lambda.00=10.\lambda.0;
\]
\[
B_4=0100=0.B_1.00=\lambda.B_2.00=0.B_3.\lambda=010.\lambda.0;
\]
\[
B_5=0000=0.\lambda.000=00.\lambda.00=000.\lambda.0.
\]
Тогда
\[
P_\Sigma=\{0,1,01,\lambda,10,100,00,000\},
\]
\[
S_\Sigma=\{0,1,00,000,\lambda\},
\]
\[
V=P_\Sigma\cap S_\Sigma\cup\{\lambda\}=\{0,1,\lambda,00,000\}.
\]
Соответствующий граф не содержит ориентированного цикла, проходящего через вершину \(\lambda\). Следовательно, код \(C(\Sigma)\) является взаимно однозначным.
\textit{(Источник: Прокопенко, разд. 3, с. 26--27.)}

\paragraph{Пример 7.}
Пусть \(A=\{a_1,a_2,a_3,a_4,a_5\}\), \(B=\{b_1,b_2,b_3\}\), а схема кодирования имеет вид
\[
\Sigma:\quad
\begin{cases}
a_1\to b_1b_2,\\
a_2\to b_1b_3b_2,\\
a_3\to b_2b_3,\\
a_4\to b_1b_2b_1b_3,\\
a_5\to b_2b_1b_2b_2b_3.
\end{cases}
\]
Сравнивая префиксы и суффиксы кодовых слов, получаем
\[
P_\Sigma\cap S_\Sigma\cup\{\lambda\}=\{b_1b_3,b_2,\lambda\}.
\]
Это и есть множество вершин \(V\) графа \(G_\Sigma\). Элементарные коды \(B_1\) и \(B_3\) не имеют нетривиальных разложений. Далее:
\[
B_2=b_1b_3b_2=b_1b_3.\lambda.b_2,
\]
\[
B_4=b_1b_2b_1b_3=\lambda.B_1.b_1b_3,
\]
\[
B_5=b_2b_1b_2b_2b_3=b_2.B_1B_3.\lambda.
\]
Граф содержит ориентированный цикл, проходящий через вершину \(\lambda\):
\[
\lambda\to b_1b_3\to b_2\to\lambda.
\]
Выписывая слова, приписанные вершинам и рёбрам этого цикла, получаем слово
\[
b_1b_2b_1b_3b_2b_1b_2b_2b_3,
\]
которое имеет два разложения по элементарным кодам:
\[
B_1B_2B_1B_3=B_4B_5.
\]
Иными словами, слова
\[
A_1=a_4a_5
\qquad\text{и}\qquad
A_2=a_1a_2a_1a_3
\]
кодируются одним и тем же словом. Следовательно, код \(C(\Sigma)\) не является свободным.
\textit{(Источник: Прокопенко, разд. 3, с. 27--28.)}

\paragraph{Итог.}
Алфавитное кодирование сопоставляет каждой букве исходного алфавита элементарный код, а слову — конкатенацию кодов его букв. Для однозначной декодируемости код должен быть разделимым; необходимым условием является неравенство Макмиллана. Проверка взаимной однозначности в общем случае сводится к построению графа \(G_\Sigma\) по нетривиальным разложениям элементарных кодов и применению теоремы Маркова: код не является свободным тогда и только тогда, когда в графе есть ориентированный цикл, проходящий через вершину \(\lambda\). Примеры 5--7 показывают, как по графу устанавливается неоднозначность или однозначность кода и как строится слово, допускающее два декодирования.
\textit{(Обобщение по: Прокопенко, разд. 2, с. 17--22; разд. 3, с. 23--28.)}

\subsubsection{Префиксные коды и их построение}

\noindent\textit{Основные источники: Прокопенко Н.Ю. Теория кодирования: учеб. пособие. Н. Новгород: ННГАСУ, 2023. Разд. 4, с. 30--38.}

\paragraph{Определение префиксного кода.}
Пусть задано конечное множество \(A=\{a_1,\dots,a_n\}\), называемое алфавитом. Последовательность букв называется словом в данном алфавите. Если \(\alpha=\alpha_1\alpha_2\), то \(\alpha_1\) называется \emph{началом} (или \emph{префиксом}) слова \(\alpha\), а \(\alpha_2\) — \emph{окончанием} (или \emph{постфиксом}) слова \(\alpha\). Если при этом \(\alpha_1\ne\lambda\) (соответственно \(\alpha_2\ne\lambda\)), то \(\alpha_1\) (соответственно \(\alpha_2\)) называется \emph{собственным началом} (соответственно \emph{собственным окончанием}) слова \(\alpha\). Иными словами, префикс (суффикс) слова называется собственным, если он отличен от пустого слова и от самого слова.

Обозначим через \(P_\alpha\) множество всех собственных префиксов слова \(\alpha\), а через \(S_\alpha\) — множество всех собственных суффиксов. Например, для \(\alpha=110\) имеем
\[
P_\alpha=\{1,11\},\qquad S_\alpha=\{0,10\}.
\]

\textbf{Определение.} Схема \(\Sigma\) (код \(C(\Sigma)\), кодирование \(F_\Sigma\)) обладает \emph{свойством префикса}, если для любых двух различных кодовых слов \(B_i\) и \(B_j\) (\(i\ne j\)) слово \(B_i\) не является префиксом слова \(B_j\). Код \(C(\Sigma)\), обладающий свойством префикса, называется \emph{префиксным кодом}.
\textit{(Источник: Прокопенко, разд. 4, с. 30.)}

Формально: для любых \(i\ne j\) и любого слова \(\beta\in B^*\) выполнено \(B_i\ne B_j\beta\). Иными словами, никакое кодовое слово не является началом другого кодового слова.

\paragraph{Связь префиксности и разделимости.}
Напомним, что код называется \emph{разделимым}, если любое слово, составленное из элементарных кодов, единственным образом разлагается на элементарные коды. Справедлива следующая теорема.

\textbf{Теорема.} Префиксная схема является разделимой.
\textit{(Источник: Прокопенко, разд. 4, с. 31.)}

\emph{Доказательство.} Предположим противное: пусть кодирование со схемой \(\Sigma\) не является разделимым. Тогда существует слово \(\beta\in F_\Sigma(A^*)\), допускающее два различных разложения на элементарные коды:
\[
\beta=\beta_{i_1}\dots\beta_{i_k}=\beta_{j_1}\dots\beta_{j_l},
\]
причём наборы индексов \((i_1,\dots,i_k)\) и \((j_1,\dots,j_l)\) различны. Рассмотрим первое место, где эти разложения отличаются. Пусть это будут кодовые слова \(\beta_{i_t}\) и \(\beta_{j_t}\). Одно из них является префиксом другого, что противоречит свойству префикса. Следовательно, предположение неверно, и префиксная схема разделима. \(\square\)

\textbf{Замечание.} Свойство быть префиксной является \emph{достаточным}, но не \emph{необходимым} условием разделимости. Например, код
\[
a\to 0,\qquad b\to 01
\]
не является префиксным, так как код буквы \(a\) является префиксом кода буквы \(b\). Однако этот код разделим: если после символа \(0\) следует \(0\), то первый символ кодирует букву \(a\), так как \(00\) не является началом никакого кодового слова; если после \(0\) следует \(1\), то пара \(01\) является кодом буквы \(b\). Таким образом, существуют разделимые, но не префиксные коды.
\textit{(Источник: Прокопенко, разд. 4, с. 31.)}

\paragraph{Представление префиксного кода с помощью дерева.}
Пусть задан префиксный код \(V=\{v_1,\dots,v_n\}\) в двоичном алфавите \(B=\{0,1\}\). Элементарные коды определяют \emph{бинарное кодовое дерево}. Из каждой вершины дерева выходит не более двух рёбер в вершины следующего яруса; левое ребро помечается символом \(0\), правое — символом \(1\). Элементарным кодам соответствуют вершины дерева, определяемые путём, идущим от корня. Если код префиксный, то элементарные коды расположены в листьях дерева.

Дерево называется \emph{насыщенным}, если из каждой вершины, не являющейся листом, в следующий ярус выходит ровно два ребра.
\textit{(Источник: Прокопенко, разд. 4.1, с. 32.)}

Для префиксных кодов и только для них множество кодовых слов совпадает с множеством концевых вершин (листьев) кодового дерева. Это свойство часто используется при построении и анализе префиксных кодов.

\paragraph{Пример 1.}
Пусть код задан таблицей:
\[
\begin{aligned}
b_1&\to 00,\\
b_2&\to 111,\\
b_3&\to 100,\\
b_4&\to 01,\\
b_5&\to 101,\\
b_6&\to 1100,\\
b_7&\to 1101.
\end{aligned}
\]
Соответствующее кодовое дерево имеет вид: из корня выходят две ветви (0 и 1). По ветви 0: из вершины 0 идут ветви 0 и 1, приводящие к \(b_1=00\) и \(b_4=01\). По ветви 1: из вершины 1 идёт ветвь 0 к \(b_3=100\) и \(b_5=101\); ветвь 1 ведёт к вершине 11, из которой идут ветви 0 (к \(b_6=1100\) и \(b_7=1101\)) и 1 (к \(b_2=111\)). Все кодовые слова расположены в листьях, дерево не является насыщенным, так как, например, из вершины 10 не выходят два ребра (она является листом для \(b_3\) и \(b_5\)? На самом деле в этом примере из вершины 10 выходят рёбра к \(b_3\) и \(b_5\), так что она не лист; но есть другие вершины, например, из вершины 0 выходят два ребра, из вершины 1 — два, из вершины 11 — два, из вершины 110 — два; все внутренние вершины имеют по два потомка, значит, дерево насыщенное). В данном примере дерево является насыщенным.
\textit{(Источник: Прокопенко, разд. 4.1, с. 32.)}

\paragraph{Построение префиксного кода по набору длин.}
Пусть задан префиксный код \(V=\{v_1,\dots,v_n\}\) в алфавите \(B=\{b_1,\dots,b_m\}\). Пусть \(|v_i|=l_i\) — длины кодовых слов. Упорядоченное множество
\[
L(V)=\langle l_1,\dots,l_n\rangle
\]
называется \emph{спектром длин} кода \(V\). Коды \(v_1\) и \(v_2\) называются \emph{спектрально эквивалентными}, если их спектры длин совпадают.

\textbf{Теорема.} Если для чисел \(l_1,\dots,l_n\) выполнено неравенство Макмиллана
\[
\frac{1}{q^{l_1}}+\frac{1}{q^{l_2}}+\dots+\frac{1}{q^{l_n}}\le 1,
\]
то существует префиксный код \(V=\{v_1,\dots,v_n\}\) в алфавите \(B=\{b_1,\dots,b_m\}\) со спектром длин \(L(V)=\langle l_1,\dots,l_n\rangle\).
\textit{(Источник: Прокопенко, разд. 4.2, с. 33--34.)}

\emph{Доказательство.} Будем полагать, что \(l_1\le l_2\le\dots\le l_n\). Разобьём множество чисел на \(k\) непересекающихся классов так, что в каждом классе стоят одинаковые числа. Пусть \(L_i\) — \(i\)-й класс, \(s_i\) — длина кодовых слов в этом классе, \(|L_i|\) — число элементов в классе. Тогда неравенство Макмиллана можно переписать в виде
\[
\frac{|L_1|}{q^{s_1}}+\frac{|L_2|}{q^{s_2}}+\dots+\frac{|L_k|}{q^{s_k}}\le 1.
\]
Это неравенство равносильно следующему:
\[
|L_i|\le q^{s_i}-|L_1|q^{s_i-s_1}-|L_2|q^{s_i-s_2}-\dots-|L_{i-1}|q^{s_i-s_{i-1}}
\]
для каждого \(i=1,\dots,k\).

Построение префиксного кода (алгоритм 1):
\begin{enumerate}[label=\arabic*)]
\item Из множества \(B^{s_1}\) выберем произвольные \(|L_1|\) слов. Обозначим их через \(V_1\subseteq B^{s_1}\). Возможность выбора следует из неравенства для \(i=1\).
\item Из множества \(B^{s_2}\) выберем произвольные \(|L_2|\) слов так, чтобы ни одно из слов \(V_2\) не являлось бы префиксом выбранных ранее. Обозначим их через \(V_2\). Выбор таких слов возможен в силу выполнимости неравенства для \(i=2\).
\item Пусть построены множества \(V_1,V_2,\dots,V_{i-1}\). На \(i\)-м шаге выберем произвольные \(|L_i|\) слов из множества
\[
B^{s_i}\big/\bigl((V_1\times B^{s_i-s_1})\cup (V_2\times B^{s_i-s_2})\cup\dots\cup (V_{i-1}\times B^{s_i-s_{i-1}})\bigr).
\]
Это множество обозначим через \(V_i\).
\item Через \(k\) шагов алгоритм заканчивается. Код \(V=\bigcup_{i=1}^k V_i\) является префиксным по построению.
\end{enumerate}
\textit{(Источник: Прокопенко, разд. 4.2, с. 34--35.)}

\paragraph{Пример 3.}
Пусть \(L(V)=\langle1,3,3,4,4,4,4\rangle\). Проверим неравенство Макмиллана:
\[
\frac{1}{2}+\frac{1}{2^3}+\frac{1}{2^3}+\frac{1}{2^4}+\frac{1}{2^4}+\frac{1}{2^4}+\frac{1}{2^4}
=\frac{1}{2}+\frac{2}{8}+\frac{4}{16}=1\le 1.
\]
Неравенство выполнено, следовательно, префиксный код с таким спектром существует.

Разобьём спектр на классы: \(L_1=\{1\}\), \(L_2=\{3,3\}\), \(L_3=\{4,4,4,4\}\). Имеем \(s_1=1\), \(s_2=3\), \(s_3=4\), \(|L_1|=1\), \(|L_2|=2\), \(|L_3|=4\).

Шаг 1. Выбираем \(V_1\subseteq\{0,1\}\), \(|V_1|=1\). Пусть \(V_1=\{1\}\).

Шаг 2. Выбираем \(V_2\subseteq\{0,1\}^3\), \(|V_2|=2\), так, чтобы слова не начинались с \(1\). Пусть \(V_2=\{000,001\}\).

Шаг 3. Выбираем \(V_3\subseteq\{0,1\}^4\), \(|V_3|=4\), так, чтобы слова не начинались с \(1\), а также не начинались с \(000\) или \(001\). Пусть \(V_3=\{0100,0101,0110,0111\}\).

Искомый код:
\[
V=V_1\cup V_2\cup V_3=\{1,\ 000,\ 001,\ 0100,\ 0101,\ 0110,\ 0111\}.
\]
Он является префиксным по построению. Заметим, что можно было построить и другой код с тем же спектром длин; все такие коды спектрально эквивалентны.
\textit{(Источник: Прокопенко, разд. 4.2, с. 35.)}

\paragraph{Пример 4.}
Для спектра \(L(V)=\langle1,1,3\rangle\) построим префиксный код в различных алфавитах.
\begin{enumerate}[label=\arabic*)]
\item В алфавите \(B=\{0,1\}\) неравенство Макмиллана не выполняется:
\[
\frac{1}{2}+\frac{1}{2}+\frac{1}{2^3}=\frac{1}{2}+\frac{1}{2}+\frac{1}{8}>1.
\]
Следовательно, префиксного кода не существует.

\item В алфавите \(B=\{0,1,2\}\) неравенство Макмиллана выполнено:
\[
\frac{1}{3}+\frac{1}{3}+\frac{1}{3^3}=\frac{19}{27}\le 1.
\]
Значит, префиксный код существует. Например:
\[
\Sigma=\{B_1=1,\ B_2=2,\ B_3=000\}
\]
или
\[
\Sigma=\{B_1=0,\ B_2=2,\ B_3=110\}.
\]
Построение по алгоритму 1: \(V_1=\{0,2\}\) (два слова длины 1), \(V_2=\{110\}\) (одно слово длины 3, не начинающееся с 0 или 2). Искомый код \(V=V_1\cup V_2=\{0,2,110\}\).

\item В алфавите \(B=\{0,1,2,3\}\) неравенство Макмиллана также выполнено:
\[
\frac{1}{4}+\frac{1}{4}+\frac{1}{4^3}=\frac{33}{64}\le 1.
\]
Например, \(\Sigma=\{B_1=0,\ B_2=2,\ B_3=100\}\) или \(\Sigma=\{B_1=2,\ B_2=3,\ B_3=000\}\).
\end{enumerate}
\textit{(Источник: Прокопенко, разд. 4.2, с. 36.)}

\paragraph{Алгоритм Шеннона построения префиксного кода.}
Пусть задан набор длин \(l_1,\dots,l_n\), удовлетворяющих неравенству Макмиллана. В силу доказанной теоремы существует префиксный код с таким спектром длин. Алгоритм Шеннона (алгоритм 2) строит этот код конструктивно.

Построим последовательность чисел \(q_1,q_2,\dots,q_n\) по правилам:
\[
q_1=0,\qquad q_{i+1}=q_i+2^{-l_i}\quad (i=1,\dots,n-1).
\]
Очевидно, \(0\le q_i<1\) и каждое \(q_i\) имеет единственное представление в виде двоичной дроби с \(l_i\) знаками после запятой:
\[
q_i=0,c_1^{(i)}c_2^{(i)}\dots c_{l_i}^{(i)}.
\]
Кодовое слово \(v_i\) определяется как последовательность этих \(l_i\) цифр:
\[
v_i=c_1^{(i)}c_2^{(i)}\dots c_{l_i}^{(i)}.
\]
Так как набор длин упорядочен по неубыванию, при \(h>i\) выполняются неравенства \(l_h\ge l_i\) и \(q_h\ge q_i+2^{-l_i}\). Поэтому кодовое слово \(v_h\) отличается от \(v_i\) в первых \(l_i\) разрядах. Следовательно, \(v_i\) не является префиксом \(v_h\), и построенный код является префиксным.
\textit{(Источник: Прокопенко, разд. 4.2, с. 37.)}

\paragraph{Пример 5.}
Пусть \(L(V)=\langle1,3,3,4,4,4,4\rangle\). Неравенство Макмиллана проверено в примере 3. Построим последовательность \(q_1,\dots,q_7\):
\[
\begin{aligned}
q_1&=0{,}0,\\
q_2&=0+2^{-1}=0{,}100,\\
q_3&=0{,}100+2^{-3}=0{,}100+0{,}001=0{,}101,\\
q_4&=0{,}101+2^{-3}=0{,}101+0{,}001=0{,}1100,\\
q_5&=0{,}1100+2^{-4}=0{,}1100+0{,}0001=0{,}1101,\\
q_6&=0{,}1101+2^{-4}=0{,}1101+0{,}0001=0{,}1110,\\
q_7&=0{,}1110+2^{-4}=0{,}1110+0{,}0001=0{,}1111.
\end{aligned}
\]
Выбирая в качестве \(B_i\) дробную часть \(q_i\), получаем код
\[
\Sigma=\{B_1=0,\ B_2=100,\ B_3=101,\ B_4=1100,\ B_5=1101,\ B_6=1110,\ B_7=1111\}.
\]
Этот код является префиксным и спектрально эквивалентен коду, построенному в примере 3 (там был код \(\{1,000,001,0100,0101,0110,0111\}\)).
\textit{(Источник: Прокопенко, разд. 4.2, с. 38.)}

\paragraph{Замечание о выборе алгоритма.}
Оба алгоритма (алгоритм 1 и алгоритм Шеннона) позволяют построить префиксный код по заданному спектру длин. Алгоритм 1 действует последовательно по классам длин и более нагляден; алгоритм Шеннона использует двоичные дроби и удобен для программирования. Результаты работы алгоритмов могут быть различными, но все построенные коды спектрально эквивалентны, то есть имеют одинаковый набор длин кодовых слов.
\textit{(Источник: Прокопенко, разд. 4.2, с. 34--38.)}

\paragraph{Итог.}
Префиксный код — это код, в котором ни одно кодовое слово не является префиксом другого. Префиксные коды являются разделимыми, и их удобно представлять с помощью кодового дерева, где кодовые слова соответствуют листьям. Основная задача — построение префиксного кода по заданному набору длин. Необходимым условием существования такого кода является неравенство Макмиллана; для префиксных кодов оно оказывается и достаточным. Алгоритм 1 строит код последовательным выбором слов из множеств \(B^{s_i}\), исключая уже занятые префиксы. Алгоритм Шеннона использует двоичные дроби \(q_i\) и выбирает первые \(l_i\) цифр их двоичного разложения. Оба алгоритма дают префиксный код с заданным спектром длин, что подтверждается примерами.
\textit{(Обобщение по: Прокопенко, разд. 4, с. 30--38.)}

\subsubsection{Оптимальное кодирование (сжатие информации)}

\noindent\textit{Основные источники: Прокопенко Н.Ю. Теория кодирования: учеб. пособие. Н. Новгород: ННГАСУ, 2023. Разд. 5, с. 39--57.}

\paragraph{Постановка задачи сжатия.}
Вопросы сжатия информации играют важную роль в информатике, что связано с развитием вычислительной техники и средств связи, а также с необходимостью хранения и передачи больших объёмов информации. Одно из направлений теории кодирования, в котором изучаются вопросы сжатия, называется \emph{экономным кодированием}.

\textbf{Определение.} Методы кодирования, которые позволяют построить (без потери информации) коды сообщений, имеющие меньшую длину по сравнению с исходным сообщением, называются \emph{методами сжатия} (или \emph{упаковки}) информации.
\textit{(Источник: Прокопенко, разд. 5, с. 39.)}

Качество сжатия обычно определяется \emph{коэффициентом сжатия}, который измеряется в процентах и показывает, насколько сжатое сообщение короче исходного. Например, коэффициент сжатия \(0{,}6\) означает, что сжатые данные занимают \(60\%\) от исходного размера. Значения, большие \(1\), говорят о том, что выходной файл больше входного (отрицательное сжатие). Также используется коэффициент сжатия в \emph{bpb} (bit per bit) — число бит сжатого файла на один бит входного. Величина \(100\cdot(1-k)\), где \(k\) — коэффициент сжатия, также отражает качество сжатия: значение \(60\) означает, что файл занимает на \(60\%\) меньше места.
\textit{(Источник: Прокопенко, разд. 5, с. 39--40.)}

Возможности сжатия определяются вероятностными и структурными (синтаксическими) свойствами информации. С помощью вероятностей моделируются частоты появления различных букв в сообщениях, частоты фрагментов, частоты самих сообщений. Дополнительные возможности появляются, когда рассматриваются не любые последовательности символов, а только некоторые (например, в текстах русского языка не встречаются фрагменты «ггг», «аь» и т.п.). Начало математическому исследованию экономного кодирования положила работа К. Шеннона «Математическая теория связи» (1948). Вопросы сжатия с учётом структурных свойств изучал А.А. Марков.
\textit{(Источник: Прокопенко, разд. 5, с. 40.)}

\paragraph{Стоимость кодирования.}
Пусть на буквах алфавита \(A=\{a_1,\dots,a_n\}\) задано распределение вероятностей \(P=\{p_1,\dots,p_n\}\), \(p_i\ge 0\), \(\sum_{i=1}^n p_i=1\), и пусть \(V=\{v_1,\dots,v_n\}\) — свободный код. \emph{Стоимостью кодирования} \(f\) называется величина
\[
C_f(P)=\sum_{i=1}^n p_i |v_i|,
\]
где \(|v_i|\) — длина элементарного кода буквы \(a_i\). Стоимость показывает, во сколько раз в среднем увеличивается длина кодируемого сообщения, то есть число букв, приходящееся на одну кодируемую букву.
\textit{(Источник: Прокопенко, разд. 5, с. 41.)}

Требование экономичности заключается в минимизации средней длины кодового слова. Положим
\[
C^*(P)=\min_f C_f(P).
\]
Код \(V^*\) со схемой \(f^*\), для которой \(C_{f^*}(P)=C^*(P)\), называется \emph{оптимальным} для набора вероятностей \(P\). Оптимальные коды дают в среднем минимальное увеличение длин слов.

\paragraph{Свойства оптимальных кодов.}
Существует несколько способов получения кодов с малой избыточностью, причём все они обладают следующими свойствами:
\begin{enumerate}[label=\arabic*)]
\item Более вероятным буквам источника соответствуют более короткие кодовые слова (принцип статистического кодирования).
\item Никакое кодовое слово не является началом другого, более длинного.
\item Все буквы алфавита, используемого для передачи по каналу, приблизительно равновероятны.
\item Символы в последовательности на выходе кодера практически независимы.
\end{enumerate}
Коды с одинаковым спектром длин элементарных кодов имеют одинаковую стоимость. Поэтому оптимальные коды ищут в классе префиксных кодов.
\textit{(Источник: Прокопенко, разд. 5, с. 42.)}

\paragraph{Пример 1.}
Пусть \(A=\{a_1,a_2,a_3,a_4\}\), \(P=\{0{,}4;0{,}25;0{,}2;0{,}15\}\). Рассмотрим две схемы кодирования:
\[
f_1:\quad
\begin{cases}
a_1\to 00,\\
a_2\to 01,\\
a_3\to 10,\\
a_4\to 11;
\end{cases}
\qquad
f_2:\quad
\begin{cases}
a_1\to 1,\\
a_2\to 01,\\
a_3\to 000,\\
a_4\to 001.
\end{cases}
\]
Для \(f_1\) стоимость равна
\[
C_{f_1}(P)=0{,}4\cdot2+0{,}25\cdot2+0{,}2\cdot2+0{,}15\cdot2=2.
\]
Для \(f_2\):
\[
C_{f_2}(P)=0{,}4\cdot1+0{,}25\cdot2+0{,}2\cdot3+0{,}15\cdot3=0{,}4+0{,}5+0{,}6+0{,}45=1{,}95.
\]
Таким образом, стоимость кодирования может изменяться при переходе от одной схемы к другой.
\textit{(Источник: Прокопенко, разд. 5, с. 41--42.)}

\paragraph{Алгоритм Фано.}
Алгоритм Фано строит код, близкий к оптимальному. Упорядоченный по невозрастанию вероятностей список букв делится на две последовательные части так, чтобы суммы вероятностей входящих в них букв как можно меньше отличались друг от друга. Буквам из первой части приписывается символ \(0\), буквам из второй — символ \(1\). Затем точно так же поступают с каждой из полученных частей, если она содержит хотя бы две буквы. Процесс продолжается до тех пор, пока весь список не разобьётся на части, содержащие по одной букве. Каждой букве ставится в соответствие последовательность символов, приписанных данной букве в результате этого процесса.
\textit{(Источник: Прокопенко, разд. 5.1, с. 42--43.)}

\textbf{Пример 2.} Для списка сообщений с заданным распределением частот построим код Фано.
\[
\begin{array}{c|cccccccc}
a_i & a_1 & a_2 & a_3 & a_4 & a_5 & a_6 & a_7 & a_8\\
\hline
p_i & 0{,}08 & 0{,}1 & 0{,}15 & 0{,}15 & 0{,}3 & 0{,}05 & 0{,}12 & 0{,}05
\end{array}
\]
Упорядочим по убыванию: \(a_5(0{,}3)\), \(a_3(0{,}15)\), \(a_4(0{,}15)\), \(a_7(0{,}12)\), \(a_2(0{,}1)\), \(a_1(0{,}08)\), \(a_6(0{,}05)\), \(a_8(0{,}05)\).
Разделим на две части: \(\{a_5,a_3\}\) (сумма \(0{,}45\)) и \(\{a_4,a_7,a_2,a_1,a_6,a_8\}\) (сумма \(0{,}55\)). Разница \(0{,}1\) минимальна. Припишем первой части \(0\), второй — \(1\).
В первой части \(\{a_5,a_3\}\): \(a_5\) получает \(00\), \(a_3\) — \(01\).
Во второй части \(\{a_4,a_7,a_2,a_1,a_6,a_8\}\): делим на \(\{a_4,a_7\}\) (\(0{,}27\)) и \(\{a_2,a_1,a_6,a_8\}\) (\(0{,}28\)). Приписываем \(0\) и \(1\).
Далее: \(a_4\to 100\), \(a_7\to 101\). Для \(\{a_2,a_1,a_6,a_8\}\): делим на \(\{a_2\}\) (\(0{,}1\)) и \(\{a_1,a_6,a_8\}\) (\(0{,}18\)). \(a_2\to 110\). Для \(\{a_1,a_6,a_8\}\): делим на \(\{a_1\}\) (\(0{,}08\)) и \(\{a_6,a_8\}\) (\(0{,}1\)). \(a_1\to 1110\), \(a_6\to 11110\), \(a_8\to 11111\).
Итоговый код:
\[
a_5\to 00,\quad a_3\to 01,\quad a_4\to 100,\quad a_7\to 101,\quad a_2\to 110,\quad a_1\to 1110,\quad a_6\to 11110,\quad a_8\to 11111.
\]
Стоимость кодирования:
\[
C_\Sigma(P)=2\cdot0{,}3+2\cdot0{,}15+3\cdot0{,}15+3\cdot0{,}12+3\cdot0{,}1+4\cdot0{,}08+5\cdot0{,}05+5\cdot0{,}05=2{,}83.
\]
\textit{(Источник: Прокопенко, разд. 5.1, с. 43--44.)}

\paragraph{Алгоритм Шеннона.}
Алгоритм Шеннона применим, когда все вероятности \(p_i>0\). Буквы упорядочиваются по невозрастанию вероятностей. Каждой букве \(a_i\) ставится в соответствие двоичная дробь \(q_i\), определяемая как сумма вероятностей предыдущих букв:
\[
q_1=0,\qquad q_{i+1}=q_i+p_i.
\]
Кодовое слово \(v_i\) — это первые \(l_i\) знаков двоичной записи \(q_i\), где \(l_i=\lceil -\log_2 p_i\rceil\). Построенный код является префиксным.
\textit{(Источник: Прокопенко, разд. 5.2, с. 45--46.)}

\textbf{Пример.} Для того же распределения, что и в примере 2, алгоритм Шеннона даёт код:
\[
a_1\to 1000,\quad a_2\to 1001,\quad a_3\to 010,\quad a_4\to 011,\quad a_5\to 00,\quad a_6\to 10110,\quad a_7\to 1010,\quad a_8\to 10111.
\]
Стоимость кодирования:
\[
C_\Sigma(P)=4\cdot(0{,}08+0{,}1)+3\cdot(0{,}15+0{,}15)+2\cdot0{,}3+5\cdot(0{,}05+0{,}05)+4\cdot0{,}12=3{,}2.
\]
\textit{(Источник: Прокопенко, разд. 5.2, с. 45--46.)}

\paragraph{Алгоритм Хаффмана.}
Алгоритм Хаффмана — адаптивный жадный алгоритм оптимального префиксного кодирования алфавита с минимальной избыточностью. Он был разработан в 1952 году Д. Хаффманом. Идея состоит в том, что зная вероятности символов, можно построить коды переменной длины: символам с большей вероятностью ставятся в соответствие более короткие коды. Коды Хаффмана обладают свойством префиксности и являются оптимальными.

Классический алгоритм Хаффмана на входе получает таблицу частот встречаемости символов. Строится дерево кодирования Хаффмана:
\begin{enumerate}[label=\arabic*)]
\item Символы входного алфавита образуют список свободных узлов. Каждый лист имеет вес, равный вероятности (или количеству вхождений) символа.
\item Выбираются два свободных узла с наименьшими весами.
\item Создаётся их родитель с весом, равным сумме весов.
\item Родитель добавляется в список свободных узлов, а два его потомка удаляются.
\item Одной дуге, выходящей из родителя, ставится в соответствие бит \(1\), другой — бит \(0\).
\item Шаги 2–5 повторяются, пока в списке не останется один узел — корень дерева.
\item Прослеживая путь от корня к каждому листу и записывая биты на дугах, получаем кодовые слова.
\end{enumerate}
\textit{(Источник: Прокопенко, разд. 5.3, с. 46--48.)}

\textbf{Пример 4.} Пусть \(A=\{a_1,\dots,a_8\}\), \(P=\{0{,}08;0{,}1;0{,}15;0{,}15;0{,}3;0{,}05;0{,}12;0{,}05\}\). Упорядочим по убыванию вероятностей. Последовательно складывая две наименьшие вероятности, получаем дерево. Итоговый код:
\[
a_5\to 00,\quad a_3\to 010,\quad a_4\to 011,\quad a_7\to 100,\quad a_2\to 101,\quad a_1\to 1100,\quad a_6\to 11010,\quad a_8\to 11011.
\]
Стоимость кодирования:
\[
C_\Sigma(P)=2\cdot0{,}3+3\cdot(0{,}15+0{,}15+0{,}12+0{,}1+0{,}08)+4\cdot(0{,}05+0{,}05)=2{,}8.
\]
\textit{(Источник: Прокопенко, разд. 5.3, с. 48--49.)}

\paragraph{Оптимальность кода Хаффмана.}
Коды Хаффмана являются оптимальными и префиксными. Справедлива теорема:
\[
H(P)\le C^*(P)\le H(P)+1,
\]
где \(H(P)=-\sum_{i=1}^n p_i\log_2 p_i\) — энтропия источника. Коэффициент эффективности кода определяется как
\[
\gamma=\frac{H(P)}{C_\Sigma(P)}.
\]
Для оптимального кода \(\gamma=1\). Избыточность равна \(1-\gamma\). При блочном кодировании, выбирая длину блока \(N\) достаточно большой, можно сделать стоимость кодирования на одну букву сколь угодно близкой к \(H(P)\).
\textit{(Источник: Прокопенко, разд. 5.3, с. 49--50, 55.)}

\paragraph{Пример 5.}
Для распределения \(P=\{0{,}4;0{,}2;0{,}1;0{,}1;0{,}1;0{,}05;0{,}05\}\) код Хаффмана имеет стоимость \(C_\Sigma(P)=2{,}5\). Энтропия \(H(P)\approx2{,}42\). Коэффициент эффективности \(\gamma\approx0{,}9885\), избыточность \(1-\gamma\approx0{,}0115\), что соответствует возможному сокращению длины кода на \(1{,}15\%\). Код близок к оптимальному, так как для оптимального двоичного кода \(H(P)=C_\Sigma(P)\) и \(\gamma=1\).
\textit{(Источник: Прокопенко, разд. 5.3, с. 49--50.)}

\paragraph{Сравнение методов.}
Метод Хаффмана производит идеальное сжатие (сжимает данные до их энтропии), если вероятности символов точно равны отрицательным степеням числа \(2\). Результаты эффективного кодирования по методу Хаффмана всегда лучше результатов кодирования по методу Шеннона–Фано. Алгоритм Фано строит код, близкий к оптимальному, но не всегда оптимальный. Алгоритм Шеннона также даёт близкий к оптимальному код, но его стоимость может быть выше, чем у Хаффмана.
\textit{(Источник: Прокопенко, разд. 5.3, с. 50.)}

\paragraph{Блочное кодирование.}
Для передачи данных можно перейти от побуквенного кодирования к кодированию «блоков», состоящих из фиксированного числа последовательных букв. Например, при кодировании двухбуквенных комбинаций алфавита \(\{А,Б\}\) с вероятностями \(P(А)=0{,}8\), \(P(Б)=0{,}2\) по методу Фано получаем:
\[
АА\to 0,\quad АБ\to 10,\quad БА\to 110,\quad ББ\to 111.
\]
Цена кода для блока \(C_\Sigma(P)=1\cdot0{,}64+2\cdot0{,}16+3\cdot(0{,}16+0{,}04)=1{,}56\). Стоимость одной буквы \(C_\Sigma(P)/2=0{,}78\).
При кодировании по Хаффману трёхбуквенных комбинаций стоимость одной буквы уменьшается до \(0{,}728\), приближаясь к энтропии \(H(P)=0{,}723\).
\textit{(Источник: Прокопенко, разд. 5.3, с. 53--55.)}

\paragraph{Итог.}
Оптимальное кодирование направлено на минимизацию средней длины кодового слова при заданном распределении вероятностей. Стоимость кодирования определяется как математическое ожидание длины кодового слова. Алгоритмы Фано, Шеннона и Хаффмана строят префиксные коды; наилучшим является алгоритм Хаффмана, который всегда даёт оптимальный код. Эффективность кода оценивается через энтропию источника и коэффициент эффективности \(\gamma=H(P)/C_\Sigma(P)\). При блочном кодировании стоимость на одну букву приближается к энтропии. Эти методы широко применяются в сжатии изображений, видео, в архиваторах и протоколах передачи данных.
\textit{(Обобщение по: Прокопенко, разд. 5, с. 39--57.)}

\subsubsection{Помехоустойчивое кодирование}

\noindent\textit{Основные источники: Прокопенко Н.Ю. Теория кодирования: учеб. пособие. Н. Новгород: ННГАСУ, 2023. Разд. 6, с. 59--92; разд. 7, с. 87--90.}

\paragraph{Назначение помехоустойчивого кодирования.}
При передаче информации по каналу связи возникают помехи, ошибки, и небольшая часть информации теряется. Помехоустойчивое кодирование необходимо для устранения ошибок, которые возникают в процессе передачи и хранения информации. Без использования помехоустойчивого кодирования было бы невозможно передавать большие объёмы информации, так как в любой системе передачи и хранения неизбежно возникают ошибки.
\textit{(Источник: Прокопенко, разд. 6, с. 59.)}

В зависимости от того, используется в системе обнаружение или исправление ошибок с помощью помехоустойчивого кода, различают следующие варианты:
\begin{itemize}
\item \emph{запрос повторной передачи} (Automatic Repeat reQuest, ARQ): с помощью помехоустойчивого кода выполняется только обнаружение ошибок, при их наличии производится запрос на повторную передачу пакета данных;
\item \emph{прямое исправление ошибок} (Forward Error Correction, FEC): производится декодирование помехоустойчивого кода, то есть исправление ошибок с его помощью.
\end{itemize}
Возможен также гибридный вариант: полученный пакет пытаются исправить, и если это не удаётся, отправляют запрос на повторную передачу.
\textit{(Источник: Прокопенко, разд. 6, с. 59.)}

Любое помехоустойчивое кодирование добавляет избыточность, за счёт чего появляется возможность восстановить информацию при частичной потере данных. В эффективном оптимальном кодировании избыточность убирали, а в помехоустойчивом кодировании добавляется контролируемая избыточность.
\textit{(Источник: Прокопенко, разд. 6, с. 59.)}

\paragraph{Классификация помехоустойчивых кодов.}
Помехоустойчивые коды делятся на \emph{непрерывные} и \emph{блочные} (блоковые). В непрерывных кодах процесс кодирования и декодирования носит непрерывный характер; сверточный код является частным случаем непрерывного кода. В блочных кодах процесс кодирования и декодирования осуществляется по блокам; с точки зрения понимания работы блочный код проще.
\textit{(Источник: Прокопенко, разд. 6.1, с. 60.)}

По используемому алфавиту коды делятся на двоичные, оперирующие битами, и недвоичные.

Блочные коды делятся на:
\begin{itemize}
\item \emph{систематические} — отдельно неизменённые информационные символы и отдельно проверочные; выходной блок состоит из информационных и проверочных символов;
\item \emph{несистематические} — символы исходного сообщения в явном виде не присутствуют; по выходному блоку нельзя понять, что было на входе.
\end{itemize}
\textit{(Источник: Прокопенко, разд. 6.1, с. 61.)}

\paragraph{Параметры помехоустойчивого кодирования.}
Первый параметр — \emph{скорость кода} \(R\), которая характеризует долю информационных («полезных») данных в сообщении:
\[
R=\frac{m}{n}=\frac{m}{m+k},
\]
где \(n\) — количество символов закодированного сообщения, \(k\) — количество проверочных символов, \(m\) — количество информационных символов. Параметры \(n\) и \(m\) часто приводят вместе с наименованием кода: например, код Хэмминга \((7,4)\) означает, что на вход кодера приходит 4 символа, а на выходе 7 символов.
\textit{(Источник: Прокопенко, разд. 6.1, с. 61--62.)}

Второй параметр — \emph{кратность обнаруживаемых ошибок}: количество ошибочных символов, которые код может обнаружить.

Третий параметр — \emph{кратность исправляемых ошибок} \(t\): количество ошибочных символов, которые код может исправить.

\emph{Избыточность кода}:
\[
R=\frac{k}{m}.
\]
Основная проблема кодирования — найти коды с максимально возможной скоростью передачи информации при минимальной избыточности.
\textit{(Источник: Прокопенко, разд. 6.1, с. 62.)}

\paragraph{Простой бит чётности.}
Самый простой метод помехоустойчивого кодирования — добавление одного бита чётности. Если в информационном сообщении нечётное количество единиц, добавляется \(0\); если чётное — добавляется \(1\). Если принятый бит чётности не совпадает с рассчитанным, считается, что произошла ошибка. Кратность исправляемых ошибок при этом равна \(0\), так как ошибки нельзя исправить, а кратность обнаруживаемых равна \(1\).
\textit{(Источник: Прокопенко, разд. 6.1, с. 62.)}

\paragraph{Прямоугольный код.}
Прямоугольный код — код с контролем чётности, позволяющий исправить одну ошибку. Последовательность из \(0\) и \(1\) составляется в прямоугольную матрицу (например, \(4\times4\)); для каждой строки и каждого столбца вычисляется бит чётности. При передаче и возникновении ошибки проверка показывает, в каком столбце и в какой строке произошла ошибка; чтобы исправить, достаточно инвертировать неверный бит. Этот код исправляет все однобитные ошибки, но не все двух- и трёхбитные. Скорость кода для прямоугольного кода \((4\times4)\): \(R=16/24=0{,}66\).
\textit{(Источник: Прокопенко, разд. 6.1, с. 63.)}

\paragraph{Кодовое расстояние.}
Важнейшим параметром кода является \emph{кодовое расстояние} \(d_\Sigma\) — минимальное расстояние Хэмминга между двумя различными кодовыми словами:
\[
d_\Sigma=\min d(x,y),\qquad x\ne y,\quad x,y\in\Sigma.
\]
\emph{Вес кодового слова} — число ненулевых символов. Для двоичного кода под расстоянием понимается расстояние Хэмминга — число позиций, в которых соответствующие символы двух кодовых слов одинаковой длины различны.
\textit{(Источник: Прокопенко, разд. 6.1, с. 63--64.)}

\textbf{Теорема 1.} Код \(\Sigma\) исправляет все комбинации из \(t\) или менее ошибок тогда и только тогда, когда кодовое расстояние не меньше \(2t+1\):
\[
d\ge 2t+1.
\]
\textbf{Теорема 2.} Код \(\Sigma\) обнаруживает все комбинации из \(t\) или менее ошибок тогда и только тогда, когда кодовое расстояние не меньше \(t+1\):
\[
d\ge t+1.
\]
\textit{(Источник: Прокопенко, разд. 6.1, с. 64.)}

Геометрическая иллюстрация: вокруг каждого кодового слова строятся сферы радиуса \(t\). Чтобы каждая ошибка кратности \(t\) была исправлена, слово с такой ошибкой должно содержаться внутри сферы, описанной только вокруг одного кодового слова; расстояние между центрами двух сфер должно быть не меньше \(2t+1\). Для обнаружения ошибки слово с ошибкой не должно быть кодовым, то есть расстояние между любыми различными кодовыми словами должно быть больше \(t\).
\textit{(Источник: Прокопенко, разд. 6.1, с. 64--65.)}

\paragraph{Пример 1.}
Дана схема кодирования:
\[
\Sigma:\quad
\begin{cases}
a_1\to 00000,\\
a_2\to 10101,\\
a_3\to 01011,\\
a_4\to 11110.
\end{cases}
\]
Матрица расстояний:
\[
\begin{array}{c|cccc}
 & a_1 & a_2 & a_3 & a_4\\
\hline
a_1 & 0 & 3 & 3 & 4\\
a_2 & 3 & 0 & 4 & 3\\
a_3 & 3 & 4 & 0 & 3\\
a_4 & 4 & 3 & 3 & 0
\end{array}
\]
Минимальное расстояние \(d_\Sigma=3\). Тогда \(3\ge t+1\), откуда \(t\le2\) для обнаружения, и \(3\ge 2t+1\), откуда \(t\le1\) для исправления. Следовательно, код позволяет обнаруживать все однократные и двукратные ошибки и исправлять любые однократные ошибки.
\textit{(Источник: Прокопенко, разд. 6.1, с. 65--66.)}

\paragraph{Код Хэмминга.}
Код Хэмминга — двоичный код, позволяющий автоматически исправлять одиночные ошибки. Пусть дано слово длины \(m\). В кодовом слове должно быть \(n=m+k\) разрядов, где \(m\) — количество информационных разрядов, \(k\) — количество контрольных разрядов. Количество контрольных разрядов выбирается так, чтобы удовлетворялось неравенство
\[
2^k-1\le n<2^k
\qquad\text{или}\qquad
m+k\le2^k-1.
\]
Минимальные значения \(k\) при заданных \(n\) приведены в таблице:

\[
\begin{array}{c|cccc}
n & 2 & 3 & 4-7 & 8-13\\
\hline
k_{\min} & 1 & 2 & 3 & 4
\end{array}
\]
\textit{(Источник: Прокопенко, разд. 6.2, с. 67.)}

Для произвольного слова \(A=a_1a_2\dots a_n\) определяется вектор
\[
H(A)=a_1e_k(1)\oplus a_2e_k(2)\oplus\dots\oplus a_ne_k(n),
\]
где \(e_k(i)\) — \(k\)-разрядная двоичная запись номера \(i\). \textbf{Теорема 3.} Код Хэмминга \(H_n\), состоящий из всех слов \(X=x_1x_2\dots x_n\), для которых \(H(X)=(00\dots0)\), является кодом с исправлением одного замещения. Число элементов кода Хэмминга \(|H_n|=2^{n-k}\), позиции с номерами \(2^0,2^1,\dots,2^{k-1}\) являются проверочными, остальные — информационными. Таким кодом можно кодировать все слова длины \(m=n-k\).
\textit{(Источник: Прокопенко, разд. 6.2, с. 68.)}

\paragraph{Пример 3.}
Для \(X=010101\) и \(Y=110100\), \(n=6\), \(k=3\).
\[
H(X)=(010)\oplus(100)\oplus(110)=(000),
\]
то есть \(X\in H_6\).
\[
H(Y)=(001)\oplus(010)\oplus(100)=(111),
\]
то есть \(Y\notin H_6\).
\textit{(Источник: Прокопенко, разд. 6.2, с. 68--69.)}

\paragraph{Пример 4.}
Пусть в слове \(X=010101\in H_6\) произошло замещение пятого символа, в результате получилось \(T=010111\). Тогда
\[
H(T)=(010)\oplus(100)\oplus(101)\oplus(110)=(101),
\]
числовое значение \(H(T)\) равно \(5\), что определяет номер замещённого символа.
\textit{(Источник: Прокопенко, разд. 6.2, с. 69.)}

\paragraph{Пример 5.}
Пусть кодируемое слово \(A=1100\), \(m=4\), \(k=3\), \(n=7\). Закодированное сообщение имеет вид \(p_1p_2p_3p_4p_5p_6p_7\), где \(p_3=1\), \(p_5=0\), \(p_6=0\), \(p_7=0\) — информационные символы, а \(p_1,p_2,p_4\) — проверочные. Проверочные символы вычисляются по формулам:
\[
p_1=p_3\oplus p_5,\qquad p_2=p_3\oplus p_6,\qquad p_4=p_5\oplus p_6.
\]
В результате \(p_1=1\), \(p_2=1\), \(p_4=0\). Закодированное слово: \(1101100\).
\textit{(Источник: Прокопенко, разд. 6.2, с. 69.)}

\paragraph{Линейно блочные коды.}
Линейные блочные коды — это класс кодов с контролем чётности, которые можно описать парой чисел \((n,m)\). В процессе кодирования блок из \(m\) символов сообщения (вектор сообщения) преобразуется в больший блок из \(n\) символов кодового слова (кодовый вектор), образованный с использованием элементов данного алфавита. Если алфавит состоит только из двух элементов (\(0\) и \(1\)), код является двоичным.

Линейным блочным \((n,m)\)-кодом называется множество \(N\) последовательностей длины \(n\), называемых кодовыми словами, которое характеризуется тем, что сумма двух кодовых слов является кодовым словом, а произведение любого кодового слова на элемент поля (скаляр) также является кодовым словом. Обычно \(N=q^m\), где \(m\) — некоторое целое число.
\textit{(Источник: Прокопенко, разд. 6.3, с. 75.)}

Линейный \((n,m)\)-код называется высокоскоростным, если отношение \(m/n\) близко к \(1\), и низкоскоростным, если \(m/n\) близко к нулю. Благодаря линейности для запоминания или перечисления всех кодовых слов достаточно хранить только те слова, которые образуют базис соответствующего линейного пространства.
\textit{(Источник: Прокопенко, разд. 6.3, с. 75.)}

\paragraph{Способы задания линейных кодов.}
\begin{enumerate}[label=\arabic*)]
\item Перечислением кодовых слов.
\item Системой проверочных уравнений:
\[
b_j=\sum_{i=1}^m a_ih_{ij},
\]
где \(j\) — номер проверочного символа, \(i\) — номер информационного символа, \(h_{ij}\) — коэффициенты, принимающие значения \(0\) или \(1\).
\item Матричный способ, основанный на построении порождающей и проверочной матриц.
\end{enumerate}
\textit{(Источник: Прокопенко, разд. 6.3, с. 76.)}

\emph{Порождающей матрицей} кода называется \(G(n,m)\)-матрица, строки которой состоят из координат векторов \(g_1,\dots,g_m\) в некотором заданном базисе. Любое кодовое слово является линейной комбинацией строк матрицы \(G\). В систематической (канонической) форме матрица имеет вид
\[
G(n,m)=[I_m,Q_{m\times k}],
\]
где \(I_m\) — единичная матрица, \(Q_{m\times k}\) — прямоугольная матрица, составленная из проверочных символов.
\textit{(Источник: Прокопенко, разд. 6.3, с. 76--77.)}

\emph{Проверочная матрица} в систематическом виде имеет вид
\[
H(n,m)=[Q_{m\times k}^T,I_k],
\]
где \(I_k\) — единичная матрица, \(Q_{m\times k}^T\) — транспонированная матрица \(Q_{m\times k}\) из порождающей матрицы. Справедливо соотношение
\[
G(n,m)\cdot H^T(n,m)=0.
\]
\textit{(Источник: Прокопенко, разд. 6.3, с. 78.)}

\paragraph{Синдром ошибки.}
\(k\)-разрядный двоичный вектор \(S\), определяемый как
\[
S=y\cdot H^T,
\]
где \(y\) — принятое слово, называется \emph{синдромом ошибки}. Если ошибок нет, то синдром равен нулю. Если произошла одиночная ошибка, то синдром равен вектору-столбцу подматрицы \(Q\), соответствующему искажённой позиции.
\textit{(Источник: Прокопенко, разд. 6.3, с. 79.)}

Основные свойства линейных кодов:
\begin{enumerate}[label=\arabic*)]
\item Произведение любого кодового слова \(v_i\) на транспонированную проверочную матрицу даёт нулевой вектор: \(v_iH^T=0\).
\item Произведение слова с ошибкой \(v'_i\) на \(H^T\) называется синдромом: \(v'_iH^T=S_i\).
\item Между порождающей и проверочной матрицами в систематическом виде существует однозначное соответствие: \(G\cdot H^T=0\).
\item Кодовое расстояние \(d_0\) равно минимальному числу линейно зависимых столбцов проверочной матрицы.
\item Произведение информационного слова на порождающую матрицу даёт кодовое слово: \(s=a\cdot G\).
\item Синдром вычисляется по правилу \(c=b\cdot H^T\).
\item Кодовое расстояние любого линейного \((n,m)\)-кода удовлетворяет неравенству \(d_0\le n-m+1\) (граница Сингтона).
\end{enumerate}
\textit{(Источник: Прокопенко, разд. 6.3, с. 79--80.)}

\paragraph{Пример 13.}
Линейный блочный \((5,2)\)-код задан производящей матрицей
\[
G=\begin{pmatrix}1&0&1&1&0\\0&1&0&1&1\end{pmatrix}.
\]
Проверочная матрица:
\[
H=\begin{pmatrix}1&0&1&0&0\\1&1&0&1&0\\0&1&0&0&1\end{pmatrix}.
\]
Принят вектор \(b=00110\), известны только одиночные ошибки. Декодирование по минимуму расстояния: расстояния до кодовых слов \((00000)\), \((10110)\), \((01011)\), \((11101)\) равны соответственно \(2\), \(1\), \(3\), \(2\). Минимальное расстояние \(1\) до слова \(10110\), следовательно, передавался вектор \(10110\), информационная последовательность \(10\).

Декодирование вычислением синдрома: \(c=bH^T=(110)\). Синдром не равен нулю, есть ошибка. Вектору \(c=110\) соответствует вектор ошибки в первом символе \(e=10000\). Суммируя \(b\oplus e=10110\), получаем тот же результат.
\textit{(Источник: Прокопенко, разд. 6.3, с. 80--82.)}

\paragraph{Пример 14.}
Построить код Хэмминга с параметрами \(d_\Sigma=3\), \(m=3\). Проверочная матрица строится из всевозможных 3-разрядных двоичных чисел, исключая нуль:
\[
H=\begin{pmatrix}0&0&0&1&1&1&1\\0&1&1&0&0&1&1\\1&0&1&0&1&0&1\end{pmatrix}.
\]
После приведения к каноническому виду и составления проверочных уравнений получаем код \((7,4)\). Принятый вектор \(b=1011001\) имеет синдром \(c=bH^T=(001)\). Синдром совпадает с первым столбцом матрицы \(H\), значит, ошибка в первом символе; передавался вектор \(a=0011001\).
\textit{(Источник: Прокопенко, разд. 6.3, с. 82--83.)}

\paragraph{Циклические коды (кратко).}
Циклическим кодом называется линейный блочный \((n,m)\)-код, который характеризуется свойством цикличности: сдвиг влево на один шаг любого разрешённого кодового слова даёт также разрешённое кодовое слово, принадлежащее этому же коду. Множество кодовых слов представляется совокупностью многочленов степени \((n-1)\) и менее:
\[
v(x)=v_{n-1}x^{n-1}+\dots+v_0x^0.
\]
Общее свойство кодовых слов циклического кода — их делимость без остатка на некоторый многочлен \(g(x)\), называемый \emph{порождающим}. Результатом деления двучлена \(x^n+1\) на \(g(x)\) является \emph{проверочный многочлен} \(h(x)\).
\textit{(Источник: Прокопенко, разд. 7.1, с. 87.)}

При декодировании циклических кодов используются \emph{многочлен ошибок} \(e(x)\) и \emph{синдромный многочлен} \(S(x)\):
\[
e(x)=v'(x)+v(x),\qquad
S(x)=R_{g(x)}[v'(x)],
\]
где \(v'(x)\) — принятое (с ошибкой) кодовое слово, \(v(x)\) — переданное. Ненулевые коэффициенты \(e(x)\) занимают позиции, соответствующие ошибкам. В процессе декодирования по принятому слову вычисляется синдром, затем по таблице находится соответствующий многочлен \(e(x)\), суммирование которого с принятым словом даёт исправленное кодовое слово.
\textit{(Источник: Прокопенко, разд. 7.1, с. 87--88.)}

Циклический код может быть задан порождающей и проверочной матрицами. Для несистематического кода матрицы строятся циклическим сдвигом порождающего и проверочного многочленов:
\[
G(n,m)=\|g(x);xg(x);\dots;x^{m-1}g(x)\|^T,
\]
\[
H(n,m)=\|h(x);xh(x);\dots;x^{k-1}h(x)\|^T.
\]
Для систематического циклического кода матрица \(G(n,m)\) имеет вид \([I_m,R_{m\times k}]\), где строки \(R_{m\times k}\) определяются из выражений \(r_i(x)=R_{g(x)}[a_i(x)x^k]\).
\textit{(Источник: Прокопенко, разд. 7.2, с. 88--90.)}

\paragraph{Пример 20.}
Для циклического \((7,4)\)-кода с порождающим многочленом \(g(x)=x^3+x+1\) матрицы \(G\) и \(H\) имеют вид:
\[
G=\begin{pmatrix}
1&0&0&0&1&0&1\\
0&1&0&0&1&1&1\\
0&0&1&0&1&1&0\\
0&0&0&1&0&1&1
\end{pmatrix},
\qquad
H=\begin{pmatrix}
1&1&1&0&1&0&0\\
0&1&1&1&0&1&0\\
1&1&0&1&0&0&1
\end{pmatrix},
\]
где \(h(x)=(x^7+1)/g(x)=x^4+x^2+x+1\).
\textit{(Источник: Прокопенко, разд. 7.2, с. 89--90.)}

\paragraph{Итог.}
Помехоустойчивое кодирование добавляет контролируемую избыточность для обнаружения и исправления ошибок. Основные параметры кода — скорость \(R=m/n\), избыточность \(k/m\), кратность обнаруживаемых и исправляемых ошибок. Минимальное кодовое расстояние \(d_\Sigma\) определяет возможности кода: для исправления \(t\) ошибок необходимо \(d\ge2t+1\), для обнаружения \(t\) ошибок — \(d\ge t+1\). Код Хэмминга исправляет одиночные ошибки; для него \(n=m+k\), \(m+k\le2^k-1\), проверочные разряды располагаются на позициях \(2^0,2^1,\dots\), а синдром указывает номер искажённого разряда. Линейно блочные \((n,m)\)-коды задаются порождающей и проверочной матрицами, связанными соотношением \(GH^T=0\); декодирование выполняется по синдрому \(S=yH^T\). Циклические коды используют аппарат многочленов: кодовые слова делятся на порождающий многочлен \(g(x)\), а синдромный многочлен позволяет локализовать ошибки.
\textit{(Обобщение по: Прокопенко, разд. 6, с. 59--92; разд. 7, с. 87--90.)}

\subsubsection{Краткое резюме}

Алфавитное кодирование сопоставляет словам исходного алфавита кодовые слова; для однозначной декодируемости необходимо выполнение неравенства Макмиллана, а достаточным условием является префиксность. Префиксные коды удобно строить по заданному спектру длин (алгоритмы 1 и Шеннона). Для сжатия информации используются оптимальные коды Фано, Шеннона и Хаффмана, минимизирующие среднюю длину кодового слова. Помехоустойчивое кодирование добавляет избыточность для обнаружения и исправления ошибок; минимальное кодовое расстояние определяет кратность обнаруживаемых и исправляемых ошибок. Код Хэмминга исправляет одиночные ошибки, линейно блочные коды задаются порождающей и проверочной матрицами и декодируются с помощью синдрома, циклические коды используют аппарат многочленов.
\textit{(Обобщение по: Прокопенко, разд. 2--7, с. 17--92.)}

\noindent\textit{Источники к пункту 1.6:}
\begin{itemize}
\item Прокопенко Н.Ю. Теория кодирования: учеб. пособие. Н. Новгород: ННГАСУ, 2023. 94 с.
\item Разд. 2 «Алфавитное кодирование», с. 17--22.
\item Разд. 3 «Алгоритм Маркова распознавания взаимной однозначности алфавитного кодирования», с. 23--28.
\item Разд. 4 «Построение префиксного кода», с. 30--38.
\item Разд. 5 «Оптимальное кодирование», с. 39--57.
\item Разд. 6 «Помехоустойчивое кодирование», с. 59--92.
\item Разд. 7 «Циклические коды», с. 87--90.
\end{itemize}

\end{document}